\label{chap:p3-83-44-clasificacion_numero}


\label{chap:cobertura-taxonomia-numeros}
\index{HMT number!taxonomy}
\index{Hensel transducer!complete displacement}
\index{rational number!periodicity}

\cref{chap:numero-medicion} defined an HMT number as a compatible section
of the inverse survival system, and not as a sequence of
digits devoid of genealogy. There remained, however, two questions that
should be separated precisely. The first is a question of
\emph{representation}: which real values can the
3-adic transducer represent before the particular closure filters are imposed? The second is
a question of \emph{selection}: which sections of the raw space are
selected by APP, TRIT, and TPK when their compatibilities are required
simultaneously? This section resolves the first, formulates the second as
an explicit domain restriction, and extracts from the base-one-thousand chart a
complete taxonomy of visible values. The foundational tools are
\(p\)-adic arithmetic, symbolic dynamics, and the positional characterization
of the rational numbers \cite{Koblitz1984,LindMarcus1995,HardyWright2008}.

The distinction is structural. A language capable of representing every tape does not, by that fact alone, select the tape of \(\pi\), \(e\), \(\varphi\), or \(\alpha\). Conversely, a selection law cannot be evaluated unless the space on which it acts has first been constructed. Representation and selection are two successive theorems, not two names for the same step.

\section{Unfiltered Hensel transducer and full shift on the ternary fibers}

Let \(p\) be a prime, let \(d\geq1\), and let
\(A\in\operatorname{GL}_d(\mathbb Z)\). For a row vector
\(x_t\in\mathbb Z_p^d\), we define
\begin{equation}
 y_t=x_tA,
 \qquad
 u_t=y_t\bmod p\in\mathbb F_p^d,
 \qquad
 x_{t+1}=\frac{y_t-u_t}{p}.
 \label{eq:transductor-hensel-general}
\end{equation}
The block \(u_t\) is the visible emission, and \(x_{t+1}\) is the memory that
remains after that digit is removed.

\begin{theorem}[Hensel Conjugation with Full Shift]
\label{thm:shift-hensel-completo}
For every \(T\geq1\), the map
\[
 \Psi_T:(\mathbb Z/p^T\mathbb Z)^d
 \longrightarrow(\mathbb F_p^d)^T,
 \qquad
 x_0\longmapsto(u_0,\ldots,u_{T-1}),
\]
is bijective. Its inverse is
\begin{equation}
 x_0\equiv
 \sum_{t=0}^{T-1}p^t u_tA^{-(t+1)}
 \pmod {p^T}.
 \label{eq:inversa-hensel-finita}
\end{equation}
Taking the inverse limit yields a homeomorphism.
\[
 \Psi:\mathbb Z_p^d\xrightarrow{\ \sim\ }
 (\mathbb F_p^d)^{\mathbb N}.
\]
If \(S_A\) denotes the update \(x_t\mapsto x_{t+1}\) and
\(\sigma\) removes the first block from a tape, then
\[
 \Psi\circ S_A=\sigma\circ\Psi.
\]
\end{theorem}

\begin{proof}
From \cref{eq:transductor-hensel-general} is obtained
\[
 x_t=(u_t+px_{t+1})A^{-1}.
\]
Iteration of this identity gives
\[
 x_0=
 \sum_{t=0}^{T-1}p^t u_tA^{-(t+1)}
 +p^Tx_TA^{-T}.
\]
The final summand vanishes modulo \(p^T\), which proves that a word of
\(T\) blocks determines at most one initial class. Conversely, given
any word, \(x_T=0\) is fixed and reconstruction proceeds backward by means of
\(x_t=(u_t+px_{t+1})A^{-1}\). The equality
\(x_tA=u_t+px_{t+1}\) verifies that the resulting class emits the prescribed
word. Therefore, \(\Psi_T\) is bijective.

The bijections are compatible with reduction modulo \(p^T\). Their inverse
limits produce the homeomorphism \(\Psi\); equivalently, the series of
\cref{eq:inversa-hensel-finita} converges in \(\mathbb Z_p^d\). Applying
\(S_A\) removes exactly the block \(u_0\), and therefore corresponds to
applying \(\sigma\) on the strip.
\end{proof}

In the HMT instance, \(p=3\), \(d=6\), and the published matrix has
determinant one. Hence
\begin{equation}
 \mathbb Z_3^6\simeq(\mathbb F_3^6)^{\mathbb N},
 \qquad |\mathbb F_3^6|=3^6=729.
 \label{eq:hmt-shift-729}
\end{equation}
Each tape of six-trit blocks has a unique initial 3-adic memory.
The matrix determines the coordinatization and the reading dynamics; the
universality of the alphabet proves that it excludes no tape.

\begin{corollary}[Entropy and Haar measure]
The raw Hensel update is conjugate to the full one-sided shift on
\(729\) symbols. Its topological entropy is \(6\log3\). The Haar measure on
\(\mathbb Z_3^6\) is transported to the uniform product measure on the
blocks.
\end{corollary}

\begin{proof}
There are \(729^T=3^{6T}\) cylinders of length \(T\), and each class modulo
\(3^T\) has Haar measure \(3^{-6T}\). Conjugation by
\cref{thm:shift-hensel-completo} transports both properties to the
shift.
\end{proof}

The probabilistic conclusion concerns the raw space endowed with Haar measure. It
does not assert that the published sections of the constants were sampled
according to that measure or that they are algorithmically random.

\section{Balanced ternary chart of the integers}

The raw 3-adic chart and the Archimedean chart do not exhaust the numerical types
of the HMT classification. For the integral operator, the empty word is used together with the
reduced words
\[
 \mathscr W_{\mathrm{bal}}
 :=\{\varnothing\}\cup
 \{a_1\cdots a_m:
 a_j\in\{-1,0,+1\},\ a_1\ne0\}.
\]
We define recursively
\[
 b(\varnothing)=0,
 \qquad
 b(wa)=3b(w)+a.
\]

\begin{theorem}[Unique balanced ternary representation]
\label{pf84:C03}
The evaluation
\[
 b:\mathscr W_{\mathrm{bal}}\longrightarrow\ZZ
\]
is bijective. Words whose first trit is \(+1\) represent the
positive integers, and those beginning with \(-1\) represent the negative integers.
\end{theorem}

\begin{proof}
For \(n\ne0\), the ultimo trit \(a\) is the representante unique of
\(n\bmod3\) in \(\{-1,0,+1\}\). The padre
\[
 \frac{n-a}{3}
\]
has strictly smaller absolute value. The iteration terminates at \(0\), so
that every integer has a finite expansion. The uniqueness of the balanced
residue at each step proves the uniqueness of the word. The sign is fixed
by the first nonzero trit.
\end{proof}

\begin{remark}
This bijection is an exact arithmetic chart. It does not identify all
balanced words with surviving TPK sections, nor does it replace the
genealogy, carry, or orientation of an enriched route.
\end{remark}

\section{Surjectivity of the unfiltered Archimedean chart onto \(\mathbb R\)}

Let us concatenate the coordinates of the blocks
\(u_t\in\mathbb F_3^6\) to obtain a tape
\((a_n)_{n\geq1}\in\{0,1,2\}^{\mathbb N}\). Its ternary evaluation is
\[
 \operatorname{ev}_3((a_n))=
 \sum_{n\geq1}\frac{a_n}{3^n}\in[0,1].
\]
We adopt the canonical convention that excludes a representation with an
eventually constant tail of twos only when the same point has another expansion
in the chart. The endpoint \(1\) therefore retains its unique fractional
representation \(0.222\ldots{}_3\).

\begin{theorem}[Archimedean raw coverage]
\label{thm:cobertura-arquimediana-bruta}
The map
\[
 P_{\rm raw}:\mathbb Z_3^6
 \xrightarrow{\Psi}(\mathbb F_3^6)^{\mathbb N}
 \xrightarrow{\operatorname{flat}}\{0,1,2\}^{\mathbb N}
 \xrightarrow{\operatorname{ev}_3}[0,1]
\]
is continuous and surjective. Therefore,
\[
 \mathbb Z_3^6/\!\sim_{P_{\rm raw}}
 \ \cong\ [0,1]
\]
as the value space, where \(x\sim_{P_{\rm raw}}y\) if their evaluations coincide.
\end{theorem}

\begin{proof}
Every ternary digit belongs to a unique consecutive block of six digits, so \(\operatorname{flat}\) is a bijection of tape spaces. Every \(r\in[0,1]\) has a ternary expansion; the \cref{thm:shift-hensel-completo} provides the unique 3-adic memory that emits its blocks. Continuity is verified by cylinders: fixing \(T\) blocks fixes \(6T\) digits and confines the value to an interval of diameter \(3^{-6T}\). Identification of the quotient with the image is the canonical factorization of a surjective map by its fibers.
\end{proof}

A single 3-adic chart is compact and therefore cannot continuously cover
an unbounded line. The correct globalization uses a countable family
of charts.

\begin{corollary}[Cover Archimedean countable of \(\mathbb R\)]
\label{cor:atlas-hmt-reales}
Let
\[
 \widehat X_{\rm raw}=\mathbb Z\times\mathbb Z_3^6,
 \qquad
 \widehat P_{\rm raw}(m,x)=m+P_{\rm raw}(x).
\]
Then \(\widehat P_{\rm raw}\) is surjective onto \(\mathbb R\).
Consequently, every real number admits at least one raw Hensel genealogy and
\[
 \widehat X_{\rm raw}/\!\sim_{\widehat P_{\rm raw}}
 \ \cong\ \mathbb R
\]
as a set of values.
\end{corollary}

\begin{proof}
For \(r\in\mathbb R\), tome \(m=\lfloor r\rfloor\) y aplique the
\cref{thm:cobertura-arquimediana-bruta} to \(r-m\in[0,1)\).
\end{proof}

This result makes precise the assertion that \(\mathbb R\) is the Archimedean
quotient image of a genealogical structure: the projection forgets 3-adic memory and
preserves the evaluated point. The assertion proved here concerns the raw completion
\emph{raw}. If \(X_{\rm surv}\subseteq\mathbb Z_3^6\) denotes the subspace
that also satisfies a concrete family of APP--TRIT--TPK closures, then
\[
 P_{\rm raw}(X_{\rm surv})=[0,1]
\]
is an additional theorem. The transducer guarantees full capacity; survival
determines what portion of that capacity the restricted system
realizes.

Nor is it deduced here that the native sum and product of TPK descend to
the entire global quotient. The usual arithmetic may be transported from
\(\mathbb R\), but proving that it coincides with operations constructed before
evaluation requires the corresponding descent theorem.

In this construction, the cylinders are not nonzero infinitesimals of
\(\mathbb R\). Every finite-depth cylinder has positive diameter,
and the sequence of diameters converges to zero. Two sections that agree up to
depth \(T\) and differ thereafter form a \emph{tail
perturbation}; their Archimedean difference is bounded by the diameter of the common
cylinder. This is the exact formulation of the ``cylinders tending to zero''
within the ordinary real line.

\section{Positional chart in base \(10^3\) and classification of the published sections}

The cubic completion
\(1000=729+243+27+1\) makes it natural to group the decimal digits into blocks
\(b_n\in\{0,\ldots,999\}\). For a tape of blocks, we define
\[
 \operatorname{ev}_{1000}(b_1,b_2,\ldots)
 =\sum_{n\geq1}\frac{b_n}{1000^n}.
\]
The canonical expansion excludes an eventually constant tail of blocks
\(999\). This convention resolves the duplication of cylinder endpoints.

\begin{theorem}[Visible taxonomy based on thousand]
\label{thm:taxonomia-base-mil}
Let \(x\in[0,1)\) be given, and let \((b_n)\) be its canonical base-one-thousand expansion.
Then:
\begin{enumerate}
\item \(x\) has a terminating expansion if and only if, when written as
an irreducible fraction \(a/q\), one has \(q\mid1000^N\) for some \(N\);
equivalently, the only primes dividing \(q\) are \(2\) and \(5\).
\item \(x\) is rational if and only if \((b_n)\) is eventually periodic.
\item \(x\) is irrational if and only if \((b_n)\) is not eventually
periodic.
\end{enumerate}
\end{theorem}

\begin{proof}
An expansion terminating at depth \(N\) has the form
\(D/1000^N\); upon reduction, its denominator contains only the primes \(2\) and
\(5\). Conversely, if \(q\mid1000^N\), then \(a/q\) is written with
\(N\) blocks.

If, after a prefix of length \(r\), a period of length
\(s\) is repeated whose base-thousand integer is \(P\), the tail is a geometric series:
\[
 \frac{P}{1000^r(1000^s-1)},
\]
and the value is rational. Conversely, when \(a/q\) is expanded, the remainders of
successive division can assume only \(q\) values. A zero remainder produces a
terminating expansion; repetition of a remainder produces a period. The third
statement is the exact negation of the second.
\end{proof}

The word \emph{fraction} denotes a presentation \(a/q\) of a rational
number, not an additional numerical species. Nor should periodicity and
algebraicity be conflated. A real number is algebraic when it is a zero of
some nonzero polynomial over \(\mathbb Q[X]\); it is transcendental otherwise. This
second classification is independent of the base used to write its digits.
For example,
\[
 \varphi^2-\varphi-1=0
\]
makes \(\varphi\) an algebraic irrational, whereas the transcendence
of \(e\) and \(\pi\) proceeds from classical theorems external to HMT\@.

\section{Concatenation, nonadic reduction, and carry}

Compatibility between the cubic chart and the digital root is an arithmetic
identity, not a visual resemblance. If \(u_j\in\{0,\ldots,999\}\) and
\[
 D_n=1000D_{n-1}+u_n,
 \qquad D_0=0,
\]
then \(D_n\) is the integer obtained by concatenating the first \(n\)
blocks.

\begin{theorem}[Nonadic congruence of control of the cubic chart]
\label{thm:checksum-nonadico-base-mil}
For every \(n\geq1\),
\[
 D_n\equiv\sum_{j=1}^n u_j\pmod9.
\]
Therefore, concatenation in base ten, reduction modulo nine, and digital root are compatible.
\end{theorem}

\begin{proof}
Since \(1000\equiv1\pmod9\), the recurrence gives
\(D_n\equiv D_{n-1}+u_n\pmod9\). Iteration from \(D_0=0\) produces the
formula.
\end{proof}

In the dodecaphase closure of \(\alpha\), the blockwise carry equation
has the form
\[
 P_m+E_m-\Phi_m-K_m+c_{m+1}=A_m+1000c_m.
\]
Reduced modulo nine,
\[
 A_m\equiv
 P_m+E_m-\Phi_m-K_m+c_{m+1}-c_m
 \pmod9.
\]
If the extreme carries vanish, the telescoping sum eliminates the intermediate
memory. The global residue of the output then coincides with the residue
of the combination of inputs. This law explains why the nonadic record
controls decimal concatenation; by itself, it does not select the blocks that
must be concatenated.

\section{Nonadic elevated coordinate and phase return}

Reduction modulo nine must not be confused with the identification of the
integers that possess the same residue. Let
\[
 I_9=\{1,2,\ldots,9\}.
\]
For each positive integer \(n\), we define its \emph{nonadic phase} and its
\emph{revolution index} by
\begin{equation}
 \rho_9(n)=1+((n-1)\bmod9),
 \qquad
 \kappa_9(n)=\frac{n-\rho_9(n)}9.
 \label{eq:coordenada-nonadica-levantada}
\end{equation}
The first coordinate is the digital root written in the alphabet
\(I_9\), so that the zero residue class is represented by phase \(9\).
The second records how many complete turns have preceded the
visible phase.

\begin{theorem}[Nonadic decomposition with memory]
\label{thm:descomposicion-nonadica-memoria}
The map
\[
 \Theta_9:\mathbb N_{>0}\longrightarrow\mathbb N_0\times I_9,
 \qquad
 n\longmapsto\bigl(\kappa_9(n),\rho_9(n)\bigr),
\]
is bijective, with inverse
\[
 \Theta_9^{-1}(q,r)=9q+r.
\]
In these coordinates, the arithmetic successor is conjugated with the dynamics.
\begin{equation}
 \mathcal S_9(q,r)=
 \begin{cases}
  (q,r+1),&1\leq r<9,\\
  (q+1,1),&r=9.
 \end{cases}
 \label{eq:sucesor-nonadico-levantado}
\end{equation}
In particular, the visible transition \(9\to1\) is a phase return accompanied
by the irreversible increment \(q\to q+1\); it is not an equality between the
integers \(9\) and \(10\).
\end{theorem}

\begin{proof}
Euclidean division of \(n-1\) by \(9\) uniquely provides
\(n-1=9q+s\), with \(q\in\mathbb N_0\) and \(0\leq s<9\). Taking
\(r=s+1\) yields \(n=9q+r\), with \(r\in I_9\), which proves simultaneously
existence, uniqueness, and the inverse formula. If \(r<9\), adding one unit only
increments \(r\). If \(r=9\), the equality
\(9q+9+1=9(q+1)+1\) gives the second branch of
\cref{eq:sucesor-nonadico-levantado}.
\end{proof}

This elevation separates three notions that the isolated figure mixes:
\begin{enumerate}
\item the unit of advancement, realized by the successor;
\item the visible phase \(r\), which traverses the nonadic cycle;
\item the accumulated memory \(q\), which distinguishes returns of the same phase.
\end{enumerate}
The projection \((q,r)\mapsto r\) forgets the number of turns and satisfies
\[
 \rho_9(n+9)=\rho_9(n),
 \qquad
 \kappa_9(n+9)=\kappa_9(n)+1.
\]
This is the elementary arithmetic model of the mechanism that, in the enriched
TPK state, additionally preserves Hensel quotients, orientation, signature, and
boundary datum. A return of the visible word can thus coexist with a
distinct genealogy.

Zero does not belong to \(I_9\): it is added as a separate base point before \(\Theta_9\) is applied. For negative integers, the positive modulus of the decomposition is preserved and an orientation \(\varepsilon\in\{-1,+1\}\) is appended. Consequently, neither the sign nor zero is deduced from the residual phase. This typing prevents a cyclic chart from being converted into a false arithmetic identity.

\begin{theorem}[Oriented Integer and Genealogical Sign]
\label{thm:entero-orientado-signo-genealogico}
The map
\[
\Theta_9^\pm:
\mathbb Z\setminus\{0\}
\longrightarrow
\{-1,+1\}\times\mathbb N_0\times I_9,
\]
\[
n\longmapsto
\bigl(\operatorname{sgn}(n),
      \kappa_9(|n|),
      \rho_9(|n|)\bigr)
\]
is bijective, with inverse
\[
(\varepsilon,q,r)\longmapsto\varepsilon(9q+r).
\]
Aritmetical negation acts by
\[
(\varepsilon,q,r)\longmapsto(-\varepsilon,q,r):
\]
reverses the orientation and preserves magnitude, phase, and turn memory.
\end{theorem}

\begin{proof}
The \cref{thm:descomposicion-nonadica-memoria} uniquely provides
\(|n|=9q+r\), with \(q\in\mathbb N_0\) and \(r\in I_9\). The sign of a nonzero
integer is likewise unique. The inverse formula recomposes \(n\), and replacing
\(n\) with \(-n\) changes only \(\varepsilon\).
\end{proof}

Thus, a negative is neither a special residual phase nor an
``impossible'' quantity. It is the same nonadic magnitude transported with the opposite
genealogical orientation. This structure will be prolonged in Dimensional Mechanics
by the particle--antiparticle anti-section.

Iterating \(\kappa_9\) produces a finite tower for every ordinary integer; in the inverse completion, a compatible sequence of quotients may be infinite. The distinction between the two cases anticipates the HMT taxonomy: the visible value is determined by the positional expansion, whereas the genealogical type is determined by the complete evolution of its memories.

\section{2 classification axes: visible value and transport genealogy}

The preceding taxonomy classifies the projected value. HMT adds a second
axis that ordinary decimal notation discards. For a section
\(x\in\Omega_\infty^{\rm Arch}\), we distinguish:
\[
 \begin{array}{rcl}
 \text{visible type}&=&
 \text{terminating, periodic or nonperiodic; algebraic or transcendental},\\
 \text{genealogical type}&=&
 \text{carry class, orientation, closure and monodromy of }x.
 \end{array}
\]
Two sections may share the same real value and differ on the second axis. A
rational real has an eventually periodic visible expansion, but a
redundant enriched presentation may preserve nonperiodic internal memory
that evaluation forgets. Periodicity of the \emph{value} always
refers to its canonical expansion; periodicity of the
\emph{genealogy} refers to the TPK enriched state.

This separation provides a concise definition:
\begin{equation}
 \boxed{
 \text{HMT number}
 =
 \text{compatible section with memory};
 \qquad
 \text{real value}
 =
 \text{its evaluation class}.}
 \label{eq:numero-hmt-y-valor-real}
\end{equation}
Measurement in this chart consists in coupling the section to the reader, fixing a cylinder, and
conditioning the compatible prolongations. Increasing precision means
narrowing the cylinder, not consuming the genealogy that remains in the fibre.

\section{Memorization of arbitrary depth in finite emission systems}

The visible recurrence of nine phases can coexist with an irrational
output only if the complete state preserves information that does not reduce to
those nine labels. The following result makes this necessity exact.

\begin{theorem}[Characterization via finite emitters]
\label{thm:emisor-finito-racional}
Let \(Q\) be a finite set, \(F:Q\to Q\) a deterministic transition,
\(o:Q\to\{0,\ldots,999\}\) an output, and \(q_0\in Q\). The tape
\(b_n=o(F^{n-1}(q_0))\) is eventually periodic and its evaluation is rational.
Conversely, every eventually periodic tape admits a finite deterministic
emitter.
\end{theorem}

\begin{proof}
Among \(|Q|+1\) consecutive states, two are equal. From the first
repetition, determinism forces the entire future to repeat, so the
tape is eventually periodic. The \cref{thm:taxonomia-base-mil} gives
rationality. For the converse, it suffices to construct a chain of states for the
prefix and a cycle for the period.
\end{proof}

Therefore, a law generating exactly \(e\), \(\pi\), or any other
irrational number cannot consist solely of a deterministic finite-state
automaton that repeats without additional memory. It may use an inverse
limit, 3-adic memory, increasing depth, or an
aperiodic input; all these mechanisms provide unbounded effective state.
In HMT, visible-phase return
\(g_{\rm vis}(n+9)=g_{\rm vis}(n)\) does not force return of the carry,
orientation, or boundary. That difference is precisely where
an aperiodic expansion may reside.

\section{The local square of \texorpdfstring{\(e\)}{e} and the memory break}

For a decimal word \(a_1\cdots a_n\), we shall call a factorization a \emph{local square}
of length \(\ell\) at position \(i\)
\[
 a_i\cdots a_{i+2\ell-1}=vv,
 \qquad |v|=\ell.
\]
This notion belongs to word combinatorics. It does not assert that the
infinite tail has a period.

\begin{theorem}[Finite characterization of the square \(1828^2\)]
\label{thm:cuadrado-local-e}
Consider the \(243\) catalog words and the fixed reader
\(w_6\mapsto w_{30}\mapsto d_1d_2d_3d_4\), with four triads in base one thousand.
There is a unique word whose reading contains a square of length four
beginning at the second decimal digit. It is
\[
 w_e=201101,
 \qquad
 718281828459=7\,\underbrace{1828\,1828}_{(1828)^2}\,459.
\]
The repetition is not prolonged to a third copy: the digit that would require that
continuation would be \(1\), whereas the effective digit is \(4\). In the whole
catalogue there exists only one other square of length four,
\[
 575157518472=\underbrace{5751\,5751}_{\text{local square}}\,8472,
\]
which begins at the first digit.
\end{theorem}

\begin{proof}
The exhaustive enumeration evaluates the \(243\) words using integer
arithmetic. The square profile for lengths \(1,\ldots,6\) contains
\[
 240,\ 21,\ 3,\ 2,\ 0,\ 0
\]
occurrences, distributed among
\[
 153,\ 19,\ 3,\ 2,\ 0,\ 0
\]
words. The two occurrences of length four are exactly those in the
statement. Two independent integer enumerations reproduce the same census,
and their complete outputs coincide term by term.
\end{proof}

The word \(w_e\) also has the lifting chain
\[
201101\mid121221\mid102011\mid012222\mid102011.
\]
The third and fifth visible blocks coincide. The states that transport them,
however, do not coincide. Before both emissions, their reductions
modulo \(27\) are
\[
(25,11,7,17,8,16),
\qquad
(7,8,4,23,26,7),
\]
and the subsequent Hensel quotients are
\[
(6,13,9,2,13,1),
\qquad
(15,0,3,19,23,25).
\]
Therefore, the return \(102011\) is a return of the visible projection, not
of the enriched state. This is the exact formulation of the local repetition
that is broken: the quotient preserves a genealogy that the residual block does not
record.

Consequently, the beginning
\[
e=2.718281828459\ldots
\]
does not exhibit a ``semiperiodicity,'' but rather a local word-square with
visible-phase return and memory advance without state periodicity. The result is exhaustive within the
fixed reader; its content is combinatorial and dynamical, not an inference of
irrationality from twelve digits.

\section[2 finite measures of the emitter 468 a 243]{2 finite measures of the \texorpdfstring{\(468\to243\)}{468→243} emitter: uniform frequency and structural priority}

The catalogue publico distinguishes \(468\) emissions y \(104\,976\) semillas.
Be
\[
 q:E_{468}\longrightarrow W_{243}
\]
the map that forgets the emission signature and preserves the six-trit
word. There are two natural measures because there are two finite counting universes:
\[
 \mu_{468}(w)=\frac{|q^{-1}(w)|}{468},
 \qquad
 \mu_{\rm seed}(w)=
 \frac{\sum_{e\in q^{-1}(w)}m(e)}{104\,976},
\]
where \(m(e)\) is the multiplicity of semillas of the emision \(e\).

\begin{theorem}[Canonical measures relative to the counted universe]
\label{thm:medidas-emisor-468}
The uniform measure on \(E_{468}\) is the unique probability invariant under
all permutations of its emissions; its image measure is \(\mu_{468}\).
The uniform measure on the \(104\,976\) seeds induces
\(\mu_{\rm seed}\). Their exact histograms are
\[
\begin{array}{c|rrrrrr}
|q^{-1}(w)|&1&2&3&4&5&6\\ \hline
\#\{w\}&132&66&18&3&6&18
\end{array}
\]
and
\[
\begin{array}{c|rrrrrrrrr}
\text{seed weight}
&144&288&432&576&864&1008&1440&1944&2088\\ \hline
\#\{w\}
&120&54&18&9&15&6&3&12&6.
\end{array}
\]
For the three words distinguished by the reader, one obtains
\[
\begin{array}{c|cc}
&\mu_{468}&\mu_{\rm seed}\\ \hline
\pi&5/468&7/729\\
e&1/468&1/729\\
\varphi&1/468&1/729.
\end{array}
\]
\end{theorem}

\begin{proof}
Invariance under the symmetric group assigns equal masses to all
points of \(E_{468}\); normalization fixes them at \(1/468\). The formula for
the pushforward measure follows by summing over each fiber. The same argument applied to the
seed set yields the second measure. The histograms are obtained
by exhaustive enumeration of the finite catalogue of \(468\) emissions and \(243\)
words; their sums verify
\[
132+66+18+3+6+18=243,
\]
\[
132+2\cdot66+3\cdot18+4\cdot3+5\cdot6+6\cdot18=468
\]
and the sum total of the weights is \(104\,976\).
\end{proof}

The nonuniformity of the pushforward measure is a finite theorem: distinct words have different numbers of realizations. Conversely, no normalizable uniform probability exists on all of \(\mathbb R\), and these measures do not intrinsically order all the reals. The word for \(\pi\) belongs to one of the six fibers of size five and weight \(1008\); it is not the unique maximum of either histogram. Its proved singularity arises from the composite structure of its five regions, from its oriented partition \(3+2\), and from its subsequent compatibilities, not from declaring an absolute probability on the continuum a priori.

\section{Selector orbital finite preceding the Archimedean reader}
\label{sec:selector-orbital-constantes}

The preceding measures reveal multiplicities, but an isolated six-trit word
still depends on an orientation. The catalog has an internal symmetry that
allows unoriented classes to be selected first and a representative to be fixed afterward.
Write a word as \(abc\mid def\) and define
\[
 r(abc\mid def)=cab\mid efd,
 \qquad
 s(abc\mid def)=def\mid abc.
\]
These permutations also act on the six coordinates of the datum
\(U_6\). In the complete catalog, they satisfy
\[
 r^3=1,
 \qquad s^2=1,
 \qquad srs=r^{-1},
\]
and hence they generate an action denoted by \(D_3^{\rm cat}\). The action uses
the fixed TRIT labeling
\[
 \Psi(U_6)=(-U_6\bmod3)
\]
and the published partition of the six coordinates into two triads.

For a word \(w\), let \(f(w)=|q^{-1}(w)|\) be the number of emissions,
\(m(w)\) its seed mass, and
\[
 \nu(w)=(n_0(w),n_1(w),n_2(w))
\]
its ordered census of trits. Also let \(d(e)\) be the additive diagonal class of
an emission and \(o(e)\in\{\mathrm{ES},\mathrm{NO}\}\) its published
orientation.

\begin{theorem}[Internal orbital selection of the three words]
\label{thm:selector-orbital-constantes}
The action \(D_3^{\rm cat}\) preserves the set of \(468\) emissions, the
projection \(q:E_{468}\to W_{243}\), and the seed multiplicity. Its action on
\(W_{243}\) has exactly \(43\) orbits: \(38\) of size six and
five of size three. The following uniqueness statements hold within that
quotient.

\begin{enumerate}
\item There exists a unique orbit \(\mathcal O_\pi\) of size six such that, for
every \(w\in\mathcal O_\pi\),
\[
 f(w)=5,
 \qquad m(w)=1008,
 \qquad
 \{m(e):e\in q^{-1}(w)\}=\{432,144,144,144,144\}.
\]
It is
\[
\mathcal O_\pi=
\{001112,010211,100121,112001,121100,211010\},
\]
and all its words have census \(\nu=(2,3,1)\).

\item Consider the orbits of size six whose words have a single
emission, of multiplicity \(144\), and whose joint diagonal support is
\(\{0,3,6\}\). There are exactly six. Among them there is a unique orbit
with the same census \((2,3,1)\) of \(\mathcal O_\pi\):
\[
\mathcal O_e=
\{011120,012110,101201,110012,120011,201101\}.
\]

\item In the same sector there exists a unique balanced orbit, that is, with
\(\nu=(2,2,2)\):
\[
\mathcal O_\varphi=
\{002112,020211,112002,121200,200121,211020\}.
\]
\end{enumerate}

Finally, the guiding clause
\begin{equation}
 \exists e\in q^{-1}(w):
 \qquad d(e)=6,
 \qquad o(e)=\mathrm{ES}
 \label{eq:gauge-orbital-constantes}
\end{equation}
select a single representative from each orbit and produces, respectively,
\[
 \boxed{010211,\qquad201101,\qquad121200.}
\]
\end{theorem}

\begin{proof}
The group relations are verified directly as identities of permutations of six
coordinates. Enumeration of the \(468\) rows
verifies that \(r\) and \(s\) send every \(U_6\) to another row of the catalogue with
the same multiplicity, and that ternary projection commutes with both actions.
Therefore, the orbits can be computed on the \(243\) words without losing the
fibers.

The exhaustive partition gives the size histogram.
\[
 38\cdot6+5\cdot3=243.
\]
Filtering those \(43\) orbits by the first predicate leaves exactly one;
filtering by singleton fiber, minimum multiplicity, and diagonal support
\(\{0,3,6\}\) leaves six. The census \((2,3,1)\) leaves a single orbit in the
first case, and the balanced census \((2,2,2)\) leaves a single orbit in the second.
Finally, the search for a lift with \(d=6\) and orientation
\(\mathrm{ES}\) leaves one word in each of the three orbits. All
filtrations are performed on structural descriptors of the catalogue before consulting the
decimal reader.
\end{proof}

The predicates of the theorem may be combined into a selection invariant.
For an orbit \(\mathcal O\), we define its \emph{orbital profile}
\[
 \Sigma_{\rm cat}(\mathcal O)=
 \bigl(
 |\mathcal O|,\ f,\ m,\ \mathcal M_{\rm fib},\
 \mathcal D,\ \nu
 \bigr),
\]
where \(f\) is the fibre size, \(m\) the seed mass,
\(\mathcal M_{\rm fib}\) the multiset of multiplicities of their
lifts, \(\mathcal D\) the diagonal support when used, and \(\nu\)
the trit census. The three distinguished classes occupy different positions
in this invariant space:
\[
\begin{array}{c|c|c|c|c|c}
\mathcal O&|\mathcal O|&f&m&\mathcal M_{\rm fib}&(\mathcal D,\nu)\\ \hline
\mathcal O_\pi&6&5&1008&\{432,144,144,144,144\}&(-,(2,3,1))\\
\mathcal O_e&6&1&144&\{144\}&(\{0,3,6\},(2,3,1))\\
\mathcal O_\varphi&6&1&144&\{144\}&(\{0,3,6\},(2,2,2)).
\end{array}
\]
The demonstrated singularity does not mean that these numbers receive an
absolute probability over \(\mathbb R\). It means that, within the finite universe
produced by the emitter and before evaluating digits, their classes are
separable by internal invariants and that the class of \(\pi\) has a
regional fiber of a different type. This is the mathematical formulation of a
structural priority in HMT\@.

The negative control is informative. If the diagonal support \(\{0,3,6\}\) is omitted, there exist four balanced singleton orbits and \(\mathcal O_\varphi\) ceases to be unique. Before the gauge is fixed, each orbit admits six possible orientations. Once it is fixed, the word of \(\pi\) preserves three positive lifts of \(U_6\), rather than a single region. Thus, the theorem selects the oriented word while respecting the geometric multiplicity of the five regions.

The non-uniformity is reinforced when passing to orbits:
\[
\begin{array}{c|cc}
&\text{emissions}&\text{seeds}\\ \hline
\mathcal O_\pi&5/78&14/243\\
\mathcal O_e&1/78&2/243\\
\mathcal O_\varphi&1/78&2/243.
\end{array}
\]
These quotients remain masses in the finite catalog, not probabilities
assigned to the evaluated real numbers.

The causal advance is precise. The three words cease to be named inputs: they are recovered from the symmetry and finite invariants of APP--TPK, relative to the action \(D_3^{\rm cat}\), the TRIT labeling, the distribution \(3+3\), and the orienting gauge \cref{eq:gauge-orbital-constantes}. The proof uses neither \(L_0\), \(L_1\), decimal digits, \(\alpha\), CODATA, nor the dodecaphasic state. The absolute uniqueness of that action among all possible automorphisms, the deduction of the gauge rather than its declaration, the canonicity of the reader \(L_0/L_1\) independently of the target value, and the coinductive emission of the infinite expansions remain subsequent claims.

\section{Position of \texorpdfstring{\(\pi,e,\varphi\) and \(\alpha\)}{pi, e, phi, and alpha} in the classification}

The \cref{thm:shift-hensel-completo} explains why any decimal prefix
can be re-emitted from a 3-adic state: the raw capacity is
universal. Therefore, the mere back-and-forth operation on a tape selects no
constant. The distinctive content of the published construction lies
in the objects preceding and following that re-emission:
\begin{enumerate}
\item the finite fibers and the five distinct regions of \(\pi\);
\item the typed Archimedean reader that associates words, cylinders, and blocks;
\item the dodecafase signed and the closure reversible
\[
U_{12}^{\rm sgn}
\longleftrightarrow K
\longleftrightarrow(D_3K,D_4K,Q)
\longleftrightarrow\alpha_{\rm HMT};
\]
\item the contrast criteria that distinguish an equivalence of charts
from a selector preceding evaluation.
\end{enumerate}

The twelve published decimal triples of \(\pi,e,\varphi\) are exact finite
coordinates of the enriched reader. The printed string of thirty-six
digits is the rational dodecaphasic projection
\(\alpha_{12}=\pi_{12}(\alpha_{\rm HMT})\):
\[
\frac{7297352569283800997285105472380663}{10^{36}}.
\]
The exact identity in this window is
\[
\alpha_{12}
:=\pi_{12}(\alpha_{\rm HMT})
=\frac{7297352569283800997285105472380663}{10^{36}}.
\]
The compatible recurrence at arbitrary depth constructs the infinite
section
\[
 \alpha_{\rm HMT}:=\varprojlim_N A^{(N)},
\]
and each rational window is an exact projection of that same object. The
kernel \(E_{21/22}\) constructs, in another chart, the finite simple positive
root \(\alpha_E\); the graded transport of the TPK extends that kernel into
a compatible family whose limit root is \(\alpha_B\). If
\(\alpha_A:=\varprojlim_N\rho_N(\alpha_{12})\), the truncation squares
of the two families commute and determine the same surviving
cylinder at each depth. The exact two-way theorem therefore establishes
\[
 \boxed{\alpha_A=\alpha_B=\alpha_{\rm HMT}}.
\]
The equality of \(\operatorname{Tr}_9(\alpha_E^{-1})\) and
\(\operatorname{Tr}_9(\alpha_{12}^{-1})\) is the first finite certificate of
that identity; it neither defines it, limits its scope nor selects the
extensions. The typed projection keeps the coinductive
section, its rational window and the analytic root distinct without separating their identity
in the limit.

The square root of \(-1\) belongs to another chart. In the
\cref{chap:algebras-cuadraticas-orden-nueve}, one constructs
\[
\mathbb F_9=\mathbb F_3[\iota]/(\iota^2+1),
\qquad \iota^2=-1,
\]
and the Witt matrix is proved to realize that quadratic structure. This is
an exact finite algebraic extension, not a visible category of real
expansions.

\section{Unfiltered Hensel chart, survival, and global boundary}

The following controls are falsifiers of the preceding finite results:
\begin{enumerate}
\item that two states distinct modulo \(3^T\) emit the same word of
\(T\) blocks, or that there is a word without a preimage;
\item failure of the explicit inverse of
\cref{eq:inversa-hensel-finita};
\item that a rational number should appear whose canonical expansion is not eventually
periodic, or an irrational number whose expansion is periodic;
\item that a finite deterministic emitter produces a tape that is
not eventually periodic;
\item that the catalog censuses do not sum to \(468\),
\(243\), and \(104\,976\), or that the measures of the three words do
not agree with the enumeration.
\end{enumerate}

The raw Hensel chart
\[
\mathbb Z_3^6\longrightarrow[0,1]
\]
is continuous and surjective; upon adjoining the integer part, its raw atlas covers \(\mathbb R\). This exact result does not imply that every raw fiber contains a representative that survives the additional filters. The published finite windows and the named readers of \(\pi,e,\varphi\) are verified within HMT, relative to their seeds, regions, orientations, and enriched rules.

\ifdefined\HMTInternalTraceability
For a filtered prefix, let \(T_u\) be the tree of compatible prolongations.
If \(T_u\) is finitely branching and has a node at every depth,
König's lemma yields an infinite branch; if, moreover, each stable level has
cardinality one, the branch is unique. This statement is conditional: the
universal nonemptiness of \(T_u\) for every prefix has not been proved. Thus, the
coverage of all fibers by filtered survivors and the well-defined descent of
native addition and multiplication to the Archimedean quotient remain an
\textsc{open programme}.
\fi




\label{chap:taxonomia-monodromia}

The positional representation of a number records a sequence of symbols,
but does not necessarily preserve the state determining its continuation.
Triadic Modular Holography distinguishes the two levels. Nonadic phase indicates
the position of a reading within the nine-step cycle; the enriched
state additionally preserves orientation, Euclidean quotient, carry
transport, and boundary condition. Return of the phase does not therefore
imply return of the state.

Let \(S\) be a space of enriched states, let
\(T:S\to S\) be a transition, and let
\[
g:S\longrightarrow\mathbb Z/9\mathbb Z
\]
the phase, with \(g(Ts)=g(s)+1\). For a basis \(b\geq2\), a reader
\(d:S\to\{0,\ldots,b-1\}\) produces
\[
x(s)=\sum_{n\geq0}\frac{d(T^ns)}{b^{n+1}}.
\]
When a value admits the two usual positional expansions, we adopt the canonical representation that does not end in a constant tail \(b-1\). This convention distinguishes an ambiguity of notation from multiplicity proper to HMT, which arises from distinct enriched sections having the same evaluation.
The nonadic return operator is the restriction
\[
M=T^9:g^{-1}(r)\longrightarrow g^{-1}(r).
\]
The difference between \(T^9s=s\) and \(g(T^9s)=g(s)\) formalizes the distinction
between periodicity and monodromy.

The transition of the enriched state of TPK may be given by a relation rather than a
function. In that case, the preceding notation applies to the shift on
the space of infinite histories or to a deterministic realization whose state
has been enlarged with the information necessary to fix the continuation; it is not
presupposed that a visible block has a unique successor.

\begin{definition}[HMT Number and Arithmetic Presentation]
An HMT number is solely a compatible section
\(\boldsymbol x\) of the inverse system of surviving states. Given a
positional reader \(L\) whose domain contains \(\boldsymbol x\), the pair
\((\boldsymbol x,L)\) is a presentation or measurement of the number. The
Archimedean value published by that presentation is the singleton intersection of the
associated cylinders when their diameters tend to zero. Two sections with the
same value are not identified before applying the Archimedean quotient.
\end{definition}

\begin{definition}[Visible and enriched periodicities]
The output is eventually periodic when there exist \(N\geq0\) and \(p>0\)
such that
\[
d(T^{n+p}s)=d(T^ns),\qquad n\geq N.
\]
The state is eventually periodic when
\[
T^{N+p}s=T^Ns.
\]
The second condition implies the first. The converse requires that the reader
separate the tails of the orbit under consideration, that is, that
\[
 \bigl(d(T^{n+k}s)\bigr)_{k\geq0}
 =\bigl(d(T^{m+k}s)\bigr)_{k\geq0}
 \quad\Longrightarrow\quad T^ns=T^ms.
\]
\end{definition}

\begin{theorem}[Classification by phase return]
\label{thm:clasificacion-retorno-fase}
Let \(s\in S\).
\begin{enumerate}
  \item If the orbit of \(s\) reaches a state whose future output is
  constantly null, then \(x(s)\) has a terminating expansion and is
  rational.
  \item If the state is eventually periodic, then the output is
  eventually periodic and \(x(s)\in\mathbb Q\).
  \item If the reader separates the tails of the orbit and the orbit of \(s\) is not
  eventually periodic, then the output is not eventually periodic
  and \(x(s)\notin\mathbb Q\).
  \item Nonperiodicity does not distinguish between algebraic irrationality and
  transcendence. This distinction requires a polynomial certificate or an
  additional transcendence theorem.
\end{enumerate}
\end{theorem}

\begin{proof}
The first two parts are the classical characterization of the rationals through terminating or eventually periodic positional expansions. In the third, eventual periodicity of the output would produce two identical tails. The separation hypothesis would then identify the corresponding states, contradicting the nonperiodicity of the orbit. The fourth part follows from the existence of algebraic and transcendental irrationals with expansions that are not eventually periodic.
\end{proof}

\begin{corollary}[Nonadic Monodromy and Aperiodicity]
Suppose that the phase returns every nine steps, that the return orbit
\[
 \{M^ks:k\geq0\},\qquad M=T^9,
\]
is not eventually periodic and the reader separates the tails. Then the
output is not eventually periodic, and the number produced is irrational. In this
case, phase periodicity coexists with an aperiodic evolution of the
enriched state.
\end{corollary}

The hypothesis cannot be replaced by a single inequality \(M(s)\ne s\), nor
by a finite number of distinct returns. Relative to the initial state,
the declared transition and reader, the current window certifies nine
phase returns with change of state; it proves finite monodromy, but not
by itself the aperiodicity of the entire orbit.

\ifdefined\HMTRepetirResultadosTaxonomicos
\section{The local square of \texorpdfstring{\(e\)}{e} and the memory break}

For a decimal word \(a_1\cdots a_n\), we shall call a factorization a \emph{local square}
of length \(\ell\) at position \(i\)
\[
 a_i\cdots a_{i+2\ell-1}=vv,
 \qquad |v|=\ell.
\]
This notion belongs to combinatorics on words and must not be confused with
a period of the infinite tail.

\begin{theorem}[Finite characterization of the square \(1828^2\)]
Consider the \(243\) words of the N33 catalog and the frozen reader
\(w_6\mapsto w_{30}\mapsto d_1d_2d_3d_4\), with four triads in base one thousand.
There is a unique word whose reading contains a square of length four
beginning at the second decimal digit. It is
\[
 w_e=201101,
 \qquad
 718281828459=7\,\underbrace{1828\,1828}_{(1828)^2}\,459.
\]
The repetition is not prolonged to a third copy: the next required digit
would be \(1\), whereas the effective digit is \(4\). In the whole catalogue there
exists only one other square of length four,
\[
 575157518472=\underbrace{5751\,5751}_{\text{word square}}\,8472,
\]
but begins at the first digit.
\end{theorem}

\begin{proof}
The exhaustive enumeration evaluates the \(243\) words by exact
integer arithmetic. The square profiles for lengths \(1,\ldots,6\) contain,
respectively,
\[
240, 21, 3, 2, 0, 0
\]
occurrences, distributed among
\[
153, 19, 3, 2, 0, 0
\]
words. The two occurrences of length four are exactly those written
in the statement. Two independent integer enumerations reproduce the same
census, and their complete outputs agree term by term.
\end{proof}

The word \(w_e\) also has the lifting chain
\[
201101\mid121221\mid102011\mid012222\mid102011.
\]
The third and fifth visible blocks coincide. The states that carry them,
however, do not coincide. Before both emissions, their reductions
modulo \(27\) are
\[
(25,11,7,17,8,16),
\qquad
(7,8,4,23,26,7),
\]
and the subsequent Hensel quotients are
\[
(6,13,9,2,13,1),
\qquad
(15,0,3,19,23,25).
\]
Therefore, the return \(102011\) is a return of the visible projection, not
of the enriched state. This is the exact formulation of the intuition of a
local repetition that breaks: the quotient remembers a genealogy that the
residual block does not preserve.

The result belongs to the enriched reader already fixed by the system
APP--TRIT--TPK\@. The finite census exhaustively characterizes the appearance of the
local square; the coinductive prolongation of the section belongs to the
nonadic connection constructed in the
\cref{chap:numero-heisenberg-d108}. The combinatorial property of the prefix and
the generation of the complete expansion are therefore two levels
compatible with one and the same genealogy.

This formulation provides the exact framework for interpreting the
decimal beginning of \(e\),
\[
e=2.718281828459\ldots.
\]
The correct expression is no longer ``semiperiodicity'', but \emph{local word
square with visible return and memory rupture}.

\section{2 finite measures of the emitter \(468\to243\): frequency and structural priority}

There is no uniform distribution on all real numbers. There is, however,
a canonical measure on the finite catalog when the set of \(104\,976\) TPK seeds\@
is taken to be uniform. For a visible word \(w\), define
\[
\mu_{\mathrm{TPK}}(w)
=\frac{\#\{\text{seeds projecting onto }w\}}{104\,976}.
\]
The measure is not uniform. The \(243\) words have weights ranging
from \(144\) to \(2088\), with nine distinct multiplicity values.
The exact histogram, writing the weight as the base and the number of words
as the exponent, is
\[
144^{120},\ 288^{54},\ 432^{18},\ 576^9,\ 864^{15},\
1008^6,\ 1440^3,\ 1944^{12},\ 2088^6.
\]

For the three distinguished words one obtains
\[
\mu_{\mathrm{TPK}}(010211)=\frac{1008}{104\,976}=\frac7{729},
\]
\[
\mu_{\mathrm{TPK}}(201101)
=\mu_{\mathrm{TPK}}(121200)
=\frac{144}{104\,976}=\frac1{729}.
\]
The word associated with \(\pi\) has five preimages \(U_6\) and weight seven
times greater than the words associated with \(e\) and \(\varphi\). This result
proves a combinatorial hierarchy in the catalogue; by itself, it neither identifies
the words with the classical constants nor establishes a probability
on \(\mathbb R\).

Nor does it by itself define a total priority: six words have weight
\(1008\), whereas \(e\) and \(\varphi\) belong to the minimal class of
weight \(144\), shared by another \(118\) words. The HMT distinction must
be expressed by a structural profile, for example
\[
\mathcal P(w)=
\bigl(m_{\rm semilla}(w),m_{U_6}(w),r_{\rm supervivencia}(w),
s_{\rm frontera}(w),s_{\rm excepcional}(w)\bigr),
\]
with each coordinate defined before applying external labels. A scalar of ``priority'' further requires fixing, likewise in advance, how those coordinates are ordered or weighted.
\fi

\section{Classification of constants by genealogical antecedent, generating operation, and codomain}

The HMT classification separates properties that the ordinary representation
brings together under a single value:
\begin{center}
\begin{tabular}{p{0.25\textwidth}p{0.64\textwidth}}
\toprule
Class & Structural condition \\
\midrule
Definitive & The output reaches a null tail. \\
Rational periodic & The output is eventually periodic. \\
Rational monodromy & The output is periodic, but the enriched state
traverses a nontrivial fibre that the visible reader forgets. \\
Irracional coinductivo & The output is not eventually periodic and the
cylinders determine a unique value. \\
Recurrente local & A finite word contains returns or squares, without this imposing periodicity on the suffix. \\
Algebraic & The value satisfies a nonzero polynomial with integral coefficients; the
route must be accompanied by that certificate. \\
Transcendent & The value satisfies no nonzero integer polynomial; the
classification requires a transcendence theorem. \\
Distinguido HMT & The section satisfies a frozen structural predicate of
fiber, boundary, survival, or symmetry; it is a class transverse to the
preceding ones. \\
\bottomrule
\end{tabular}
\end{center}

Thus, HMT does not alter the correct arithmetic definitions of rationality,
algebraicity, and transcendence. It refines them by adding the genealogy of the state,
phase monodromy, and the multiplicity of the fibers that Archimedean
evaluation had forgotten.

In particular, \(\varphi\) is algebraic because it satisfies
\(X^2-X-1=0\), whereas the transcendence of \(e\) and \(\pi\) follows from
classical external theorems. These properties are compatible with the three
sections being distinguished by HMT\@. The dodecaphasic closure later constructs
the exact rational projection
\(\alpha_{12}=\pi_{12}(\alpha_{\mathrm{HMT}})\), and the compatible prolongation
constructs the coinductive section \(\alpha_{\mathrm{HMT}}\), in
\cref{chap:alpha-rev7}; neither is used in this taxonomy as an input
or published here as a second residence. The kernel \(E_{21/22}\)
produces a distinct analytic root, whose agreement with the integer route is
proved in the codomain of \(\operatorname{Tr}_9\). Section, rational
projection, analytic root, reader and value thus retain their own types.

Finally, an ``infinitesimal direction'' requires its own typing. For a finite reader \(r_m\), two distinct sections in the same fiber of \(r_m\) define a difference that is invisible at depth \(m\), although active in a subsequent prolongation. They do not thereby constitute a nonzero infinitesimal real. If two sections coincide under every faithful truncation of the inverse system, they are the same section. A genuine arithmetic infinitesimal would require the construction of an ordered or valued non-Archimedean extension and proof that the HMT operations descend to it.



\section[Number as a compatible section]{Number as a compatible
section: arithmetic classification, Archimedean evaluation, and
genealogy}
\label{p3-83-c44-s1:sec:ontologia-numeros-rev11}
\index{number!value and genealogy}
\index{numerical taxonomy}

The numerical classification of HMT preserves the ordinary arithmetic species and adds a coordinate that positional evaluation discards: the genealogy of the section that publishes the value. The result is a joint classification. The first axis answers which number is obtained; the second records by which compatible sequence of states, returns, carries, and closures it is obtained. This distinction makes it possible to formulate within a single language integers, rational expansions, irrationals, the imaginary unit, regularizations, and spectral constants without identifying objects that only share a scalar.

\subsection{Archimedean evaluation separates the section from its value}

Let
\[
 X_{\infty}^{\rm surv}
 =\varprojlim_n X_n^{\rm surv}
\]
an inverse system of surviving states with its
restriction maps. A compatible section is a family
\(\boldsymbol x=(x_n)_{n\geq0}\) satisfying
\(\rho_{n+1,n}(x_{n+1})=x_n\) at every depth. A positional reader
\(L\) assigns a cylinder \(I_n(\boldsymbol x;L)\) to each prefix. When the
cylinders are compatible, their diameters converge to zero and the
convention for endpoints is fixed canonically, one defines
\begin{equation}
 \operatorname{Val}_{L}(\boldsymbol x)
 :=\text{the point of }
 \bigcap_{n\geq0}\overline{I_n(\boldsymbol x;L)}.
 \label{p3-83-c44-s1:eq:valor-lector-seccion-rev11}
\end{equation}

\HMTClaim{HMT-R-NUMBER-SECTION-VALUE}
\begin{definition}[Number, presentation and value in HMT]
\label{p3-83-c44-s1:def:numero-presentacion-valor-rev11}
An \emph{HMT number} is a compatible section
\(\boldsymbol x\in X_{\infty}^{\rm surv}\). The pair
\((\boldsymbol x,L)\) is a \emph{presentation} of the number, and
\(\operatorname{Val}_{L}(\boldsymbol x)\) is the value published by that
presentation. The fibre
\[
 \operatorname{Val}_{L}^{-1}(r)
\]
preserves the genealogies that the Archimedean quotient represents by
the same value \(r\).
\end{definition}

The capacity of the chart and selection of a
section are distinct typed results. The \cref{thm:cobertura-arquimediana-bruta} constructs
a continuous surjective map
\[
 P_{\rm raw}:\mathbb Z_3^6\longrightarrow[0,1],
\]
and the \cref{cor:atlas-hmt-reales} globalizes it to
\[
 \widehat P_{\rm raw}:\mathbb Z\times\mathbb Z_3^6
 \longrightarrow\mathbb R.
\]
These maps prove the complete representational capacity of the raw chart.
The survival predicate and the readers that distinguish
\(\pi,e,\varphi\) select, within that capacity, the corresponding structural genealogies.

\subsection{Termination and periodicity determine rationality}

For the expansion canonica in basis mil
\[
 x=\sum_{n\geq1}\frac{b_n}{1000^n},
 \qquad b_n\in\{0,\ldots,999\},
\]
one adopts the convention that eliminates the duplication produced by an
eventually constant tail of blocks \(999\). The
\cref{thm:taxonomia-base-mil} then provides the following exact classification.

\HMTClaim{HMT-R-NUMBER-DECIMAL-CLASSIFICATION}
\begin{theorem}[Visible classification in the cubic decimal chart]
\label{p3-83-c44-s1:thm:clasificacion-visible-rev11}
Let \(x=a/q\in[0,1)\) be an irreducible fraction when \(x\) is rational.
Then:
\begin{enumerate}
\item the expansion terminates if and only if \(q\mid1000^N\) for some
      \(N\), an equivalence restricting the prime factors of \(q\) to
      \(2\) and \(5\);
\item the expansion is eventually periodic if and only if
      \(x\in\mathbb Q\);
\item the expansion is not eventually periodic if and only if
      \(x\in\mathbb R\setminus\mathbb Q\).
\end{enumerate}
Algebraicity and transcendence constitute a second axis: a real number is
algebraic if it is a zero of a nonzero polynomial in \(\mathbb Q[X]\), and is
transcendental if it is a zero of none.
\end{theorem}

\begin{proof}
An expansion of length \(N\) has a denominator dividing \(1000^N\), and
the converse is obtained by expressing \(a/q\) with that denominator. For a
rational number, successive division has a finite set of remainders: a zero
remainder produces termination, and repetition of a remainder produces periodicity.
A periodic tail is, conversely, a rational geometric series. The
polynomial classification is independent of the positional base.
\end{proof}

The enriched dynamics also distinguishes output periodicity from
state periodicity. If \(T\) transports the state and \(d\) publishes a
block, the conditions
\[
 d(T^{n+p}s)=d(T^n s),
 \qquad
 T^{n+p}s=T^n s
\]
define, respectively, visible-phase return and state return. The second
implies the first. The inverse implication requires the reader to separate the
tails of the orbit, according to
\cref{thm:clasificacion-retorno-fase}. Therefore, a periodic nonadic phase nonadic periodic
can transport a nonperiodic memory periodic.

\begin{longtable}{@{}>{\raggedright\arraybackslash}p{.20\textwidth}
                    >{\raggedright\arraybackslash}p{.30\textwidth}
                    >{\raggedright\arraybackslash}p{.42\textwidth}@{}}
\caption{Joint taxonomy of the published value and the enriched state.}
\label{p3-83-c44-s1:tab:taxonomia-conjunta-numeros-rev11}\\
\toprule
\textbf{Class} & \textbf{Condition on the value} &
\textbf{Genealogical condition or required proof}\\
\midrule
\endfirsthead
\toprule
\textbf{Class} & \textbf{Condition on the value} &
\textbf{Genealogical condition or required proof}\\
\midrule
\endhead
Integer & Element of \(\mathbb Z\). & Finite and unique balanced ternary word; sign and magnitude remain typed.\\
Decimal finito & A rational number whose reduced denominator divides some power
of \(1000\). & The output reaches a null tail.\\
Rational periodic & Expansion that is eventually periodic. & An eventually periodic state suffices; a periodic output may preserve an
additional internal fibre.\\
Rational monodromic & Same rational value as the previous row. & The output returns while the enriched state traverses a fiber that the visible reader forgets.\\
Irracional coinductivo & Expansion that is not eventually periodic. & An infinite family
of compatible cylinders fixes a unique value; global aperiodicity
requires a proof over the entire orbit, not a finite window.\\
Algebraic & Root of a polynomial no zero of \(\mathbb Q[X]\). & The polynomial is appended
and its vanishing is verified.\\
Transcendent & Real external to the closure algebraic of \(\mathbb Q\). & A transcendence theorem is
appended; decimal nonperiodicity by itself
only proves irrationality.\\
Distinguido HMT & Any of the preceding arithmetic classes. & The section satisfies a frozen fiber, boundary, survival,
symmetry, or closure predicate.\\
\bottomrule
\end{longtable}

This table locates on different axes two questions that must be answered
separately. Ordinary classification determines the properties of the value;
the HMT reader determines the structural chain generating the section. For
example,
\[
 \varphi^2-\varphi-1=0
\]
classifies \(\varphi\) as an algebraic irrational, whereas the
transcendence of \(e\) and \(\pi\) belongs to external classical theorems.
All three may simultaneously be sections distinguished by different HMT
predicates. Their printed truncations of one thousand or ten thousand digits are
finite rationals that demonstrate the calculated prefixes of a prolongable output; they do not
replace the infinite object they represent.

The dodecaphase closure defines, in turn, the finite rational projection
\begin{equation}
 \alpha_{12}=\pi_{12}(\alpha_{\rm HMT})
 =\frac{7297352569283800997285105472380663}{10^{36}}.
 \label{p3-83-c44-s1:eq:alpha-racional-taxonomia-rev11}
\end{equation}
This projection is exact and rational. The object
\(\alpha_{\rm HMT}=\varprojlim_NA^{(N)}\) is the infinite section constructed
by the compatible recurrence; its rational truncations do not replace it.
The structural derivation of the limit and the classification of each
finite projection are complementary statements.

\subsection{The balanced ternary chart represents the integers
bijectively}

Let
\[
 \mathscr W_{\rm bal}
 =\{\varnothing\}\cup
 \{a_1\cdots a_m:
 a_j\in\{-1,0,+1\},\ a_1\neq0\}
\]
and define
\[
 b(\varnothing)=0,
 \qquad
 b(wa)=3b(w)+a.
\]
The \cref{pf84:C03} demonstrates that
\begin{equation}
 b:\mathscr W_{\rm bal}\xrightarrow{\ \sim\ }\mathbb Z
 \label{p3-83-c44-s1:eq:carta-entera-balanceada-rev11}
\end{equation}
is a bijection. The final trit is the unique representative of the
class modulo three, and the quotient \((n-a)/3\) strictly reduces the absolute
value until zero is reached. The first nonzero trit fixes the sign. This chart
is a representation that is exact and arithmetic; the membership of a word in a
surviving section additionally incorporates the conditions of route,
orientation, and boundary.

For \(n>0\), the coordinates
\[
 \rho_9(n)=1+((n-1)\bmod9),
 \qquad
 \kappa_9(n)=\frac{n-\rho_9(n)}9
\]
realize the decomposition \(n=9\kappa_9(n)+\rho_9(n)\). The phase
\(\rho_9\in\{1,\ldots,9\}\) publishes the zero residue class as \(9\), and
\(\kappa_9\) preserves the number of turns. With the sign
\(\varepsilon\in\{-1,+1\}\),
\[
 n\longleftrightarrow
 \bigl(\varepsilon,\kappa_9(|n|),\rho_9(|n|)\bigr)
\]
is bijective on \(\mathbb Z\setminus\{0\}\); zero remains the base
point. The visible-phase return \(9\to1\) is thus a phase transition with
memory advance, not an identification of consecutive integers.

\subsection{Finite complex structure induced by the Witt endomorphism
\texorpdfstring{\(J^2=-I_6\)}{J²=-I₆}}

The ordinary imaginary unit is the element
\(i\in\mathbb C\) satisfying \(i^2=-1\); it is algebraic of degree two over
\(\mathbb R\) and over \(\mathbb Q\). The HMT realization belongs to a
different finite domain. If \(C_P\) is the Paley conference matrix and
\(J=A_W\) its reduction modulo three in
\(\operatorname{End}_{\mathbb F_3}(\mathbb F_3^6)\), then
\begin{equation}
 \boxed{J^2=2I_6=-I_6.}
 \label{p3-83-c44-s1:eq:raiz-interna-menos-uno-ontologia-rev11}
\end{equation}
Since \(X^2+1\) is irreducible over \(\mathbb F_3[X]\), the extension
\[
 \mathbb F_9=\mathbb F_3[\iota]/(\iota^2+1)
\]
separates the eigenspaces of \(J\) associated with \(\pm\iota\). The result,
proved in \cref{sec:raiz-interna-menos-uno-rev6}, endows the ternary module with a
finite complex structure and organizes two conjugate orientations. Its force
is \emph{internally exact, proved using the declared Paley--Witt matrix}; the
comparison with \(\mathbb C\) preserves the distinction between the two fields.

\subsection{\(6\) operators realize the rational
\texorpdfstring{\( -1/12 \)}{-1/12} on distinct domains}

Zeta regularization assigns to the sequence of positive integers the finite
part obtained by meromorphic continuation:
\begin{equation}
 \sum_{n\geq1}^{[\zeta]}n
 :=\zeta(-1)=-\frac1{12}.
 \label{p3-83-c44-s1:eq:regularizacion-zeta-menos-doce-rev11}
\end{equation}
The superscript \([\zeta]\) forms part of the type of the expression. The ordinary
partial sum satisfies
\(\sum_{n=1}^{N}n=N(N+1)/2\) and grows without bound; therefore
\eqref{p3-83-c44-s1:eq:regularizacion-zeta-menos-doce-rev11} is a regularized value, not
the limit of those partial sums.

HMT additionally contains six exact algebraic realizations of
the same rational number. Scalar coincidence permits their comparison; each
object retains its domain, map, and normalization.

\begin{longtable}{@{}>{\raggedright\arraybackslash}p{.23\textwidth}
                    >{\raggedright\arraybackslash}p{.43\textwidth}
                    >{\raggedright\arraybackslash}p{.26\textwidth}@{}}
\caption{Typed realizations of \(-1/12\).}
\label{p3-83-c44-s1:tab:realizaciones-menos-doce-rev11}\\
\toprule
\textbf{Domain} & \textbf{Map and evaluation} &
\textbf{Conclusion and normalization}\\
\midrule
\endfirsthead
\toprule
\textbf{Domain} & \textbf{Map and evaluation} &
\textbf{Conclusion and normalization}\\
\midrule
\endhead
Marked hexad--octad incidence &
\(F_6\hookrightarrow F_8\), with
\(\lvert F_6\rvert=90\), \(\lvert F_8\rvert=120\), and
\[
 (90-120)/(24\binom62)=-1/12.
\]
& Exact combinatorial defect for the declared normalization.\\
Dodecaphase periodicities &
To step of \(30^\circ\) produce
\[
 \begin{aligned}
 \frac{30}{120}-\frac{30}{90}
 &=\frac14-\frac13,\\
 &=-\frac1{12}.
 \end{aligned}
\]
& Exact oriented difference of periods.\\
Gram Operator &
For \(K_{60,81}=\operatorname{diag}(0,12)\), the pseudoinverse on the
nonzero range is \(1/12\); the negative orientation publishes \(-1/12\). &
Exact operator identity on the non-vanishing component.\\
Second scale difference &
The calculo produce \(C(L)=8/81\) and the extractor declarado
\(R(L)=-(27/32)C(L)=-1/12\). & Exact after fixing the normalization factor
\(-27/32\); that factor is an explicit primitive.\\
Modular quotient eta &
\[
 \begin{aligned}
 \frac{\eta(\tau)}{\eta(3\tau)}
 &=q^{-1/12}\\[-1mm]
 &\quad{}\times
 \prod_{n\geq1}(1+q^n+q^{2n})^{-1}.
 \end{aligned}
\]
& exponent modular exact, procedente of \(1/24-3/24\).\\
Projectors of the twelve-phase cycle &
For \(D_r=I-S^r\),
\[
 \begin{aligned}
 \tau(\Pi_{\ker D_3}-\Pi_{\ker D_4})
 &=\frac{3-4}{12},\\
 &=-\frac1{12}.
 \end{aligned}
\]
& Exact transverse defect between steps three and four.\\
\bottomrule
\end{longtable}

The six realizations are proved in
\cref{sec:menos-doce-realizaciones-rev6}. Zeta regularization, the incidence defect,
the pseudoinverse, the scale extractor, the eta quotient, and the
trace of projectors are different maps. An identification of their
domains would require an additional intertwiner. Likewise, the term
\(+1/12\) of the Glaisher--Kinkelin formula belongs to the regularization
of the hyperfactorial and retains that distinct type.

\subsection{Multiplicative decomposition of primes as irreducible modes of the TPK}
\index{prime!irreducible}
\index{prime clock}

An ordinary prime is an integer \(p>1\) whose factorizations
\(p=ab\), with \(a,b\in\mathbb N\), necessarily contain a unit. The
reduced coproduct
\[
 \Delta_{\times}^{\rm red}e_n
 =\sum_{\substack{ab=n\\a,b>1}}e_a\otimes e_b
\]
turns that definition into a spectral selector. According to
\cref{pf84:C07}, the positive operator
\[
 N_{\times}^{\rm red}
 =(\overline{\Delta_{\times}^{\rm red}})^*
   \overline{\Delta_{\times}^{\rm red}}
\]
satisfies \(N_{\times}^{\rm red}e_n=d_{\rm red}(n)e_n\), and its kernel, after removing \(e_1\), selects exactly the primes. The multiplicative intertwiner transports the same selector to the native product of APP\@. This selector is the spectral projection onto \(\ker N_{\times}^{\rm red}\ominus\mathbb C e_1\), which coincides exactly with the subspace generated by the primes and does not depend on recognizing decimal lists.

Unique factorization additionally provides the Fock realization
\[
 U_F:\Gamma_s(\ell^2(\mathbb P))
 \xrightarrow{\ \sim\ }\ell^2(\mathbb N_{\geq1}),
 \qquad
 U_F\,d\Gamma(H_{\mathbb P})\,U_F^*=\log Q,
\]
where \(H_{\mathbb P}e_p=(\log p)e_p\). Each prime is a
multiplicatively irreducible mode; a composite integer is a finite occupation
of those modes.

For \(n\) coprime to \(10\), its decimal clock is the multiplicative order
\[
 \tau_n=\operatorname{ord}_n(10).
\]
If \(p\) is primo y \(p\notin\{2,3,5\}\), the
\cref{thm:cierre-nonadico-reloj-primo-rev6} demonstrates
\begin{equation}
 \tau_p\mid p-1,
 \qquad
 M_p:=\frac{10^{\tau_p}-1}{p}\in\mathbb N,
 \qquad
 M_p\equiv0\pmod9.
 \label{p3-83-c44-s1:eq:reloj-primo-cierre-nonadico-rev11}
\end{equation}
For a compound,
\[
 \operatorname{ord}_{n}(10)
 =\operatorname{lcm}_{p^a\parallel n}\operatorname{ord}_{p^a}(10).
\]
Therefore, the composite clock synchronizes prime-power clocks, whereas a prime supplies an irreducible component. The nonadic congruence \eqref{p3-83-c44-s1:eq:reloj-primo-cierre-nonadico-rev11} describes phase closure; the primality selector resides in the reduced coproduct. Separation of these two functions prevents selective capacity from being attributed to the decimal period by itself.

\subsection{Each constant is associated with its extraction operator}

The cycles of order \(3^m\), the limit of their heat traces, and the
theta--Mellin transform construct the functional
\begin{equation}
 \mathcal Z_3(s)
 =\frac{\displaystyle
 \int_0^\infty\Theta_+(x)x^{s/2-1}\,\mathrm dx}
 {\pi^{-s/2}\Gamma(s/2)}
 =\zeta(s),
 \qquad \Re s>1.
 \label{p3-83-c44-s1:eq:funcional-triadico-resumen-rev11}
\end{equation}
Theta inversion meromorphically extends the functional. The residual
partition
\[
 \zeta(s)=Z_0(s)+Z_1(s)+Z_2(s)
\]
separates aggregate, scale, and orientation. A single operatorial source admits
several extraction maps---finite part, Laurent coefficient,
evaluation, regularized derivative, or character---; coincidence of the source
does not identify the resulting constants.

% REMATE_LOCAL_CONSTANTES_TABLA
\par

\begingroup
\small
\fontsize{10pt}{11.2pt}\selectfont
\setlength{\LTpre}{0.5\baselineskip}
\setlength{\LTpost}{0.5\baselineskip}
\begin{longtable}{@{}>{\raggedright\arraybackslash}p{.15\textwidth}
                    >{\raggedright\arraybackslash}p{.245\textwidth}
                    >{\raggedright\arraybackslash}p{.305\textwidth}
                    >{\raggedright\arraybackslash}p{.195\textwidth}@{}}
\caption{Definition, extraction map, and mathematical scope of the
complementary numerical constants.}
\label{p3-83-c44-s1:tab:constantes-operadores-ontologia-rev11}\\
\toprule
\textbf{Object} & \textbf{Standard definition} &
\textbf{Map and mathematical reference} & \textbf{Conclusion and condition}\\
\midrule
\endfirsthead
\toprule
\textbf{Object} & \textbf{Standard definition} &
\textbf{Map and mathematical reference} & \textbf{Conclusion and condition}\\
\midrule
\endhead
Euler--Mascheroni \(\gamma\) &
For \(H_N=\sum_{n\leq N}n^{-1}\),
\(\gamma=\lim_{N\to\infty}(H_N-\log N)\), equivalente to the parte finite of
\(\zeta\) in \(s=1\). &
If \(C_r=\operatorname*{FP}_{s=1}Z_r(s)\), then
\(\gamma=C_0+C_1+C_2\), by
\eqref{eq:residual-three-covectors}. &
Exact identity of the residual functional demonstrated with the declared
operator and normalization.\\
Stieltjes constants \(\gamma_n\) &
If \(\zeta(1+z)-z^{-1}=\sum_{n\geq0}a_nz^n\), they are defined by
\(\gamma_n=(-1)^n n!a_n\). &
The residual coefficients \(c_{r,n}\) satisfy \(\gamma_n=(-1)^n n!\sum_r c_{r,n}\) and
\(L^{(n)}(1,\chi_{-3})/n!=c_{1,n}-c_{2,n}\): this is the coefficientwise decomposition of
\eqref{eq:stieltjes} by the three residue classes of \eqref{eq:residual-three-covectors}. &
Exact decomposition; numerical extractions verify the
coefficients, not their definition.\\
Ap\'ery \(\zeta(3)\) &
Value of the Riemann zeta function at \(s=3\). &
\(\mathcal Z_3(3)=\zeta(3)\), with prefixes bounded by rational intervals;
its spectral residence within the graded tower is fixed at
\cref{eq:bendersky,eq:odd-zeta-bendersky}. &
Exact evaluation of the functional; its irrationality follows from the
classic theorem of Apery \cite{Apery1979}.\\
Glaisher--Kinkelin \(A_{\rm GK}\) &
Regularized Growth Constant of the Hyperfactorial. &
\(A_{\rm GK}=\exp(1/12-\mathcal Z_3'(-1))\), by
\cref{eq:bendersky,eq:odd-zeta-bendersky}. &
Exact regularized derivative; \(1/12\) belongs to this hyperfactorial
normalization.\\
Euler--Kronecker constant \(\gamma_{F_{-3}}\) &
Logarithmic finite part of the Dedekind zeta function of the field
\(F_{-3}=\mathbb Q(\sqrt{-3})\). &
The residual character \(\chi_{-3}\) selects the field and produces
\(\gamma_{F_{-3}}=\gamma+L'(1,\chi_{-3})/L(1,\chi_{-3})\), through
\eqref{eq:euler-kronecker}. &
Exact identity; the choice of field comes from the character modulo three.\\
Meissel--Mertens \(B_1\) &
For \(S_P(x)=\sum_{p\leq x}p^{-1}\),
\(B_1=\lim_{x\to\infty}(S_P(x)-\log\log x)\). &
It is the finite part of the truncated prime trace over the one-particle
sector of the Fock Hamiltonian, by \eqref{eq:meissel-mertens}; the selector
of prime modes is fixed before regularization is performed. &
Exact operatorial representation once the prime modes have been fixed; the
selector of irreducibles precedes it.\\
Feigenbaum \((\delta,\alpha_{\rm F})\) &
Universal scaling constants of the renormalization fixed point for the
quadratic period-doubling cascade. &
The family dodecafasica declarada produce to level eight
\(\delta^{\mathscr B}_8(0)=4.66921218915\ldots\) and
\(\alpha^{\mathscr B}_8(0)=2.50296847988\ldots\), by
\cref{sec:p3-final:cierre-feigenbaum}. &
Exact finite calculation up to \(n=8\); convergence to the universal fixed point
requires proving convergence of the renormalization operator.
\(\alpha_{\rm F}\) and \(\alpha_{\rm HMT}\) are
distinct objects.\\
Catalan \(G\) &
\(G=L(2,\chi_{-4})=\sum_{n\geq0}(-1)^n/(2n+1)^2\). &
The natural reader is the odd character modulo four evaluated at \(2\),
by \cref{hmtctx:catalan:sec:combinaciones-lectores-vecinos-rev6}. &
Exact classical identity with localized operator; it is not an output of the
residual functional modulo three.\\
Khinchin \(K\) &
Almost-sure geometric mean of the digits of the continued fraction for the
Gauss system. &
The reader is the ergodic mean of \(\log a_1\) for
\(T(x)=x^{-1}-\lfloor x^{-1}\rfloor\) and its invariant measure, by
\cref{hmtctx:khinchin:sec:combinaciones-lectores-vecinos-rev6}. &
Incorporated classical ergodic theorem; it requires dynamics distinct from
the triadic spectral tower.\\
\bottomrule
\end{longtable}
\endgroup

The table establishes a reading rule: a constant is defined by its object, its operator, and its extraction map, not by the appearance of a decimal prefix in a catalogue. Euler--Mascheroni and Stieltjes belong to finite parts and Laurent jets; Apery, to an evaluation; Glaisher--Kinkelin, to a regularized derivative; Euler--Kronecker, to the character of the field; Meissel--Mertens, to a prime trace; Feigenbaum, to a finite dynamical family; Catalan and Khinchin, to distinct neighboring operators.

\subsection{Domains, premises, and necessary conditions of the genealogical classification of the number}

The arithmetic definitions of integer, rational, irrational, algebraic, and
transcendental are standard. The balanced ternary chart, the separation between
section and value, the native selector of irreducibles, the Paley--Witt root of
minus one, the six realizations of \(-1/12\), and the spectral readers are
deduced from the maps cited in each subsection. This section gathers them without
introducing an additional constant or numerical selector.

The force of each assertion is determined by its map:
\begin{enumerate}
\item the indicated bijections, matrix identities, coproducts, traces,
and evaluations are exact on their declared domains;
\item the transcendence of \(e\) and \(\pi\), the irrationality of
      \(\zeta(3)\) and Khinchin's ergodic theorem are classical theorems
      incorporated for purposes of control or classification;
\item the Feigenbaum approximants establish the computed branch and
      depth; universality requires the renormalization operator and
      its convergence theorem;
\item equality of two scalars from different domains authorizes a
comparison and requires an intertwiner to identify the complete
objects;
\item a decimal truncation proves the computed prefix, whereas the
      arithmetic class of the infinite value requires, as appropriate, a
      polynomial equation or a theorem of irrationality or transcendence.
\end{enumerate}

Under these rules, the HMT ontology of number acquires a precise formulation: the Archimedean value is the published coordinate; the compatible section preserves the chain that generates it; the operator determines which property is read; and its hypotheses determine exactly which conclusion follows.


\section[Genealogical closure of the concept of number]{Genealogical closure of the concept of number: compatibility, memory, and Archimedean evaluation}
\label{p3-83-c44-s2:chap:cierre-genealogico-numero}
\index{number!genealogical closure}
\index{digit!visible coordinate}
\index{value!Archimedean evaluation}
\index{evolutionary conservation of unity}

The number theory developed in HMT culminates in a distinction of types that governs the subsequent transition to Dimensional Mechanics. A \emph{digit} is a visible coordinate of a finite state; a \emph{value} is the Archimedean image published by a reader; an \emph{HMT number} is the complete compatible genealogy that preserves, together with its prefixes, orientation, quotient, carry, boundary, and memory. The fiber of a value gathers the distinct genealogies that publish that same scalar.

This construction extends ordinary arithmetic classification. Integer,
rational, irrational, algebraic, and transcendental continue to denote
properties of the value. HMT incorporates a second axis: the manner in which a
section is prolonged, returns, preserves memory, and reaches a stable
evaluation. The closure of Parts~I--III fixes both axes and the map relating them.

\subsection{Digit, section, value and symmetry}

Let
\[
 \mathcal S_0^{\rm surv}
 \xleftarrow{\rho_{1,0}}
 \mathcal S_1^{\rm surv}
 \xleftarrow{\rho_{2,1}}
 \cdots
\]
the inverse system of surviving states, where each restriction
preserves the typed TPK data belonging to the preceding level, and let
\[
 \mathcal S_\infty^{\rm surv}
 :=\varprojlim_n\mathcal S_n^{\rm surv}.
\]

\begin{definition}[Ciphers]
\label{p3-83-c44-s2:def:cifra-visible-genealogica}
Having fixed a positional reader \(L\), the digit or block visible
at depth \(n\) is
\[
 d_{L,n}(x_n):=\pi^{\rm vis}_{L,n}(x_n).
\]
It is the coordinate published by the finite presentation
\(x_n\); the fiber of \(\pi^{\rm vis}_{L,n}\) preserves the states sharing that
coordinate.
\end{definition}

\begin{definition}[HMT Number and Presentation]
\label{p3-83-c44-s2:def:numero-hmt-cierre-genealogico}
An HMT number is to compatible section
\[
 \boldsymbol x=(x_n)_{n\geq0}\in\mathcal S_\infty^{\rm surv},
 \qquad
 \rho_{n+1,n}(x_{n+1})=x_n.
\]
The pair \((\boldsymbol x,L)\) is a presentation of the number. The section
preserves the complete history; the reader declares which coordinates of that
history are published.
\end{definition}

\begin{definition}[Archimedean value]
\label{p3-83-c44-s2:def:valor-arquimediano-cierre-genealogico}
Suppose that \(L\) assigns to each \(x_n\) a compact cylinder
\(I_{L,n}(\boldsymbol x)\subset\mathbb R\), with
\[
 I_{L,n+1}(\boldsymbol x)\subseteq I_{L,n}(\boldsymbol x),
 \qquad
 \operatorname{diam}I_{L,n}(\boldsymbol x)\longrightarrow0.
\]
The value published by the presentation is
\[
 \operatorname{Val}_L(\boldsymbol x)
 :=\text{the unique element of }
 \bigcap_{n\geq0}I_{L,n}(\boldsymbol x).
\]
\end{definition}

\begin{proposition}[Existence and uniqueness of the evaluation]
\label{p3-83-c44-s2:prop:evaluacion-seccion-genealogica}
The preceding hypotheses determine a map
\[
 \operatorname{Val}_L:\mathcal S_\infty^{\rm surv}\longrightarrow\mathbb R.
\]
Each fibre \(\operatorname{Val}_L^{-1}(r)\) preserves the genealogies that the
Archimedean evaluation represents by the same value \(r\).
\end{proposition}

\begin{proof}
The finite intersection property of the nested family of compact sets
guarantees a nonempty intersection. Two distinct points would have a
positive distance bounded by every diameter, contradicting the
convergence of the latter to zero.
\end{proof}

\begin{definition}[Genealogical symmetry]
\label{p3-83-c44-s2:def:simetria-genealogica-hmt}
A symmetry of the HMT number system is a family of bijections
\(g=(g_n)_{n\geq0}\) that preserves survival and types, and satisfies
\[
 \rho_{n+1,n}\circ g_{n+1}
 =g_n\circ\rho_{n+1,n}
 \qquad(n\geq0).
\]
The action
\[
 g_\infty(\boldsymbol x):=(g_n x_n)_{n\geq0}
\]
is an automorphism of the inverse limit. A genealogical invariant is a
function constant on its orbits.
\end{definition}

Symmetry is thus characterized as a transformation compatible at every
depth. This notion renders decimal publication, exceptional
incidence, and subsequent spectral realization comparable while preserving the
proper type of each output.

\subsection{Evolutionary Conservation of Unity}

For \(n\geq1\), let
\begin{equation}
 \rho_9(n)=1+((n-1)\bmod 9),
 \qquad
 \kappa_9(n)=\frac{n-\rho_9(n)}9.
 \label{p3-83-c44-s2:eq:residuo-cociente-conservacion-unidad}
\end{equation}
Then
\[
 n=9\kappa_9(n)+\rho_9(n),
 \qquad
 \rho_9(n)\in\{1,\ldots,9\}.
\]
The coordinate \(\rho_9\) publishes the nonadic phase and represents the null class
through \(9\); \(\kappa_9\) records the number of circuits. The digital root
publishes the phase. Reduction modulo nine accompanied by its quotient publishes
phase and memory.

The associated helical realization is
\begin{equation}
 \mathcal H_9(n):=
 \left(
 \cos\frac{2\pi(\rho_9(n)-1)}9,
 \sin\frac{2\pi(\rho_9(n)-1)}9,
 \kappa_9(n)+\frac{\rho_9(n)-1}{9}
 \right).
 \label{p3-83-c44-s2:eq:helice-nonadica-conservacion-unidad}
\end{equation}

\begin{proposition}[Phase return and memory advance]
\label{p3-83-c44-s2:prop:retorno-fase-avance-memoria}
The realization \eqref{p3-83-c44-s2:eq:helice-nonadica-conservacion-unidad} satisfies
\[
 \mathcal H_9(n+9)=\mathcal H_9(n)+(0,0,1).
\]
The visible return \(9\to1\) closes the phase and increments the state
memory by one.
\end{proposition}

\begin{proof}
The identities
\(\rho_9(n+9)=\rho_9(n)\) and
\(\kappa_9(n+9)=\kappa_9(n)+1\) fix the circular coordinates and
increase the height exactly.
\end{proof}

This identity formulates evolutionary conservation: unity returns as
phase, and the new state transports the preceding closure as memory. The
cubic completion realizes the same law in three dimensions:
\begin{equation}
 10^3=(9+1)^3
 =9^3+3\cdot9^2+3\cdot9+1
 =729+243+27+1
 =729+271.
 \label{p3-83-c44-s2:eq:emrh-conservacion-evolutiva-unidad}
\end{equation}
The interior has \(729\) states; \(243+27\) constitutes the punctured
corona; the apex \(+1\) completes the base-one-thousand unit. The identity
\(729+270=999\) describes the state immediately preceding closure, and
\(729+271=1000\) describes the completion. The law
\[
 10^d=(9+1)^d
 =\sum_{r=0}^{d}\binom dr9^{d-r}
\]
transports the same unit between dimensions and preserves the genealogy of its
strata.

\subsection{Visible syntax and semantics of prolongation}

A local signature describes the current form of a branch. Its survival
semantics is the system of prolongations it can support. For
\(x_n\in\mathcal S_n^{\rm surv}\), define
\[
 \operatorname{Ext}_m(x_n)
 :=\{x_m\in\mathcal S_m^{\rm surv}:\rho_{m,n}(x_m)=x_n\},
 \qquad m\geq n.
\]
Two states that satisfy
\[
 d_{L,n}(x_n)=d_{L,n}(y_n)
\]
may possess different extension trees. In the finite witness
\(60\to81\), two branches with the same visible signature possess, respectively,
zero and twelve prolongations. The witness establishes that future continuity
adds information to the local output. The HMT number incorporates that continuity
as part of the object.

Coinduction is expressed through the compatibility of all prefixes.
Blocks of one thousand or ten thousand digits are truncated publications of a
uniformly prolongable rule. Numerical generation reaches any
requested depth; the lift to the enriched state also preserves
the hypotheses of existence, uniqueness, and compatibility of its fibers.

\subsection{Arithmetic classification and genealogical classification}

In a basis \(b\geq2\), the canonical presentation of a real number
separates the following visible classes:
\begin{enumerate}
\item an integer belongs to \(\mathbb Z\) and admits a finite balanced-ternary word;
\item a rational number has a finite or eventually periodic expansion;
\item an irrational number has an expansion that is not eventually periodic;
\item an algebraic number is a root of a nonzero polynomial in \(\mathbb Q[X]\);
\item a transcendental number belongs to the complement of the algebraic closure of
      \(\mathbb Q\) in \(\mathbb C\).
\end{enumerate}
Periodicity distinguishes rational from irrational numbers; the polynomial equation distinguishes algebraic from transcendental numbers. These are different axes. For example, \(\varphi\) is algebraic irrational because \(\varphi^2-\varphi-1=0\); the transcendence of \(e\) and \(\pi\) follows from classical theorems; an HMT genealogy additionally distinguishes the propagation, closure, and self-scaling functions that realize those values.

The same separation governs \(e+\pi\). The sum of the two presentations
defines a composite object and publishes
\[
 \operatorname{Val}(\boldsymbol e)+
 \operatorname{Val}(\boldsymbol\pi)=e+\pi.
\]
Its arithmetic class is a question in its own right. The HMT construction
determines the map that forms the composite and preserves the two genealogies
involved.

\subsection{The Euler operator family as a closure grammar}

The classical Euler equation belongs to a unique quadratic family. Let
\(J_\sigma\) be a generator satisfying
\[
 J_\sigma^2=\sigma I,
 \qquad \sigma\in\{-1,0,+1\}.
\]
The separation of even and odd powers provides
\begin{equation}
 \exp(tJ_\sigma)=C_\sigma(t)I+S_\sigma(t)J_\sigma,
 \label{p3-83-c44-s2:eq:euler-trigeometrico-cierre-numero}
\end{equation}
with the exact realizations
\[
\begin{array}{ccl}
J_{-1}^2=-I&:&\exp(tJ_{-1})=\cos t\,I+\sin t\,J_{-1},\\[1mm]
J_{0}^2=0&:&\exp(tJ_{0})=I+tJ_{0},\\[1mm]
J_{+1}^2=I&:&\exp(tJ_{+1})=\cosh t\,I+\sinh t\,J_{+1}.
\end{array}
\]
Elliptic, parabolic, and hyperbolic are three regimes of a single
operator exponential. The tritic classification determines the geometry of
the evolution.

The hyperbolic branch contains the Fibonacci self-scale. For
\[
 F=\begin{pmatrix}0&1\\1&1\end{pmatrix},
 \qquad
 A=F^2=\begin{pmatrix}1&1\\1&2\end{pmatrix},
\]
we have \(\det A=1\), \(\operatorname{tr}A=3\), and
\[
 \eta=\operatorname{arcosh}\frac32=2\log\varphi.
\]
Therefore, there exists \(J_h^2=I\) such that
\begin{equation}
 A=\exp(\eta J_h),
 \qquad
 A^n=\cosh(n\eta)I+\sinh(n\eta)J_h.
 \label{p3-83-c44-s2:eq:fibonacci-euler-hiperbolico-cierre-numero}
\end{equation}
The relation \(\varphi^2=\varphi+1\) expresses two readings of the same process:
local additive recurrence and spectral multiplicative scaling. This
compatibility provides the elementary model of the double
sum--product evaluation that APP formalizes on their common support.

\subsection{Typed Unification of Circular, Complex, and Hyperbolic Geometries}
\label{p3-83-c44-s2:sec:unificacion-compleja-circular-hiperbolica-numero}

The initial geometric motivation of HMT can be formulated exactly as a
common architecture of support, boundary, and generators. The complex plane
\(\mathbb C\) contains the unit circle
\[
 S^1=\{e^{\mathrm i\theta}:\theta\in\mathbb R\},
\]
which constitutes the elliptic realization of
\eqref{p3-83-c44-s2:eq:euler-trigeometrico-cierre-numero}. The upper half-plane
\(\mathbb H=\{w\in\mathbb C:\operatorname{Im}w>0\}\) is transported to the hyperbolic
disc by means of the Cayley transformation
\[
 \mathcal C(w)=\frac{w-\mathrm i}{w+\mathrm i},
 \qquad
 \mathcal C:\mathbb H\xrightarrow{\sim}\mathbb D,
 \qquad
 \mathbb D=\{z\in\mathbb C:|z|<1\}.
\]
The ideal boundary of \(\mathbb D\) is the same circle \(S^1\), whereas its
interior carries the Poincare metric
\[
 \mathrm ds_{\mathbb D}^2
 =\frac{4\,|\mathrm dz|^2}{(1-|z|^2)^2}.
\]
Fusion therefore consists of a typed diagram: the trigonometric
parameterization publishes the circular boundary; the complex structure
provides the chart; the Poincare metric realizes the hyperbolic interior.
The three regimes \(J^2=-I,0,+I\) also incorporate the parabolic threshold and
place this geometric concurrence within a single operator family.

\begin{proposition}[Euclidean and Hyperbolic Compatibility of the Quarter Turn]
\label{p3-83-c44-s2:prop:cuarto-giro-poincare-rev3}
Over the disco unitary
\[
 \mathbb D=\{z\in\mathbb C:|z|<1\},
\]
the application
\[
 \mathcal R_{\pi/2}:\mathbb D\longrightarrow\mathbb D,
 \qquad
 \mathcal R_{\pi/2}(z)=\mathrm i z,
\]
is simultaneously a Euclidean rotation and an isometry of the Poincaré
metric
\[
 \mathrm ds_{\mathbb D}^{2}
 =\frac{4\,|\mathrm dz|^2}{(1-|z|^2)^2}.
\]
Its extension to the boundary \(S^1=\partial\mathbb D\) is the oriented
quarter-turn. In the real realization of
\cref{thm:dos-realizaciones-fuente-cuadratica-rev3}, this map
corresponds to
\(\operatorname{Exp}((\pi/2)J_{\mathrm e})=J_{\mathrm e}\).
\end{proposition}

\begin{proof}
Multiplication by \(\mathrm i\) preserves \(|z|\) and
\(|\mathrm dz|\). It therefore preserves both the Euclidean metric and the
numerator and denominator of \(\mathrm ds_{\mathbb D}^{2}\). The exponential
identity follows from \(J_{\mathrm e}^2=-I\).
\end{proof}

The complex plane, trigonometric circle, and hyperbolic disk are not
identified as metric spaces: they share a coordinated carrier and the
quarter-turn automorphism once each realization has been declared.

\subsection{Boundary interface between quadratic rotation and minimal torsion}
\label{p3-83-c44-s2:sec:interfaz-frontera-rotacion-torsion-minima-rev3}

The function of this interface is to transport an already constructed interior
rotation to a boundary datum without confusing the quadratic endomorphism
with a scalar defect. Fix a base point
\(x_0\in\partial X\), a real bundle
\(E_{\partial}\to\partial X\), a connection \(\nabla_{\partial}\), a
boundary reader \(q_{\partial}\), and a discrete cotetrad
\(e_{\partial}\in\Omega^1(\partial X;E_{\partial})\). We write
\[
 \operatorname{Tors}(\partial X;E_{\partial})
 :=\Omega^2(\partial X;E_{\partial}),
 \qquad
 \operatorname{Hol}_{x_0}(\nabla_{\partial})
 \subseteq\operatorname{Aut}(E_{\partial,x_0}).
\]
An edge realization of the interface should construct the maps
\[
 \operatorname{Res}_{\partial}:
 \langle J_{\mathrm e}\rangle
 \longrightarrow\operatorname{Hol}_{x_0}(\nabla_{\partial}),
 \qquad
 \operatorname{Def}_{\partial}:
 \operatorname{Hol}_{x_0}(\nabla_{\partial})
 \longrightarrow\operatorname{Tors}(\partial X;E_{\partial}).
\]
The torsion determined by the connection and the co-tetrad has the type
\begin{equation}
 T_{\partial}
 :=d_{\nabla_{\partial}}e_{\partial}
 \in\Omega^2(\partial X;E_{\partial}).
 \label{p3-83-c44-s2:eq:torsion-borde-interfaz-euler-rev3}
\end{equation}
The physical realization of minimal torsion requires in addition a functional
\begin{equation}
 \Lambda_{\partial}:
 \operatorname{Tors}(\partial X;E_{\partial})\longrightarrow\mathbb R
 \label{p3-83-c44-s2:eq:funcional-borde-delta4-interfaz-rev3}
\end{equation}
such that the three maps satisfy the entanglement conditions
\begin{equation}
 \operatorname{Def}_{\partial}
 \bigl(\operatorname{Res}_{\partial}(J_{\mathrm e})\bigr)=T_{\partial},
 \qquad
 \Lambda_{\partial}(T_{\partial})=\Delta_4.
 \label{p3-83-c44-s2:eq:condiciones-borde-delta4-interfaz-rev3}
\end{equation}
These equations specify the comparator that must be constructed; they do not assert
that the three maps exist. The scalar \(\Delta_4\) is defined by
\cref{eq:delta4-electron-rev6}.

\begin{remark}[Exact condition of the boundary comparator]
The operator root \(J_{\mathrm e}^2=-I\), the quarter-turn
\(\operatorname{Exp}((\pi/2)J_{\mathrm e})=J_{\mathrm e}\), and the scalar
\(\Delta_4\) are deduced in their respective domains. The preceding boundary
chain fixes the type of the morphisms that must link them. Until
\(\operatorname{Res}_{\partial}\), \(\operatorname{Def}_{\partial}\), and
\(\Lambda_{\partial}\) are constructed, the conditions
\eqref{p3-83-c44-s2:eq:condiciones-borde-delta4-interfaz-rev3} define the required comparator
but do not imply a physical realization; the angular proximity of
its endpoints does not replace that operator.
\end{remark}

APP, TRIT, and TPK supply the discrete support of that architecture. On
\(X_9=(\mathbb Z/9\mathbb Z)^2\), the sum and product evaluations induce
projectors \(P_{\Sigma,d}\) and \(P_{\Pi,d}\) that satisfy
\[
 \|[P_{\Sigma,2},P_{\Pi,2}]\|=\frac{\sqrt2}{3},
 \qquad
 \|[P_{\Sigma,3},P_{\Pi,3}]\|=\frac{\sqrt3}{4}.
\]
This calculated incompatibility is the finite precursor of Heisenberg: two
families of observables share a support and possess incompatible spectral
decompositions. The alternating form of the torus,
\[
 \sigma((a,b),(c,d))=ad-bc\pmod9,
\]
is integrated into the Weyl--Heisenberg nonadic representation
\[
 X|n\rangle=|n+1\rangle,\qquad
 Z|n\rangle=\omega^n|n\rangle,\qquad
 ZX=\omega XZ,
 \qquad \omega=e^{2\pi\mathrm i/9}.
\]
The commutator publishes the oriented phase, and the resulting central group has
\(9^3=729\) elements.

The central block of APP simultaneously provides the local lorentzian realization:
\[
 B_c=
 \begin{pmatrix}7&2\\2&7\end{pmatrix}
 =9P_++5P_-,
 \qquad
 M_c:=\frac{B_c}{\sqrt{45}}
 =\exp\!\left(\frac12\log\frac95\,R\right).
\]
Con \(\eta_{1,1}=\operatorname{diag}(1,-1)\),
\[
 M_c^{\mathsf T}\eta_{1,1}M_c=\eta_{1,1},
 \qquad
 \det M_c=1,
 \qquad
 M_c\in SO^+(1,1).
\]
Thus Lorentz geometry and the Minkowski plane \(1+1\) emerge as an
exact spectral realization of the central block. Every subsequent dimensional
lift must intertwine this local block with its new codomain. The sequence
\[
 \boxed{
 \text{APP--TRIT--TPK}
 \longrightarrow
 \begin{cases}
 \text{elliptic--parabolic--hyperbolic family},\\
 \text{Lorentz--Minkowski block }1+1,\\
 \text{finite Weyl--Heisenberg representation}
 \end{cases}}
\]
presents the geometric and operatorial developments preceding the genealogical closure of the number.

\subsection{Constant Generation Operators: Domains, Seeds, and Evaluations}

The function of a constant is fixed through four data: mathematical
definition, domain, operator or reader, and deduced conclusion. The
\cref{p3-83-c44-s2:tab:semillas-cierre-genealogico-numero} brings together the objects that must
remain visible when number theory is closed.

\begingroup
\small
\setlength{\LTpre}{0.5\baselineskip}
\setlength{\LTpost}{0.5\baselineskip}
\begin{longtable}{@{}>{\raggedright\arraybackslash}p{.145\textwidth}
                    >{\raggedright\arraybackslash}p{.235\textwidth}
                    >{\raggedright\arraybackslash}p{.34\textwidth}
                    >{\raggedright\arraybackslash}p{.20\textwidth}@{}}
\caption{Constants and Seeds in the Genealogical Closure of the Number.}
\label{p3-83-c44-s2:tab:semillas-cierre-genealogico-numero}\\
\toprule
\textbf{Object} & \textbf{Definition or identity} &
\textbf{HMT domain and map} & \textbf{Conclusion and condition}\\
\midrule
\endfirsthead
\toprule
\textbf{Object} & \textbf{Definition or identity} &
\textbf{HMT domain and map} & \textbf{Conclusion and condition}\\
\midrule
\endhead
\(\pi,e,\varphi\) &
Sections published by compatible numerical cylinders. &
The filtration
\(w_6\to w_{30}\to R_{36}\to G_9\) produces the terminal state; the readers
distinguish closure, propagation, and self-scaling. &
For each prescribed finite depth, the map produces compatible unit
cylinders; the infinite section follows from the coinductive theorem.\\
\(\alpha_{\rm HMT}\) &
Dodecaphasic closure with carry over the terminal state. &
The double reading leads to
\((P,E,\Phi,K,c_{13}=0)\) and to the rational output of \(36\) digits; the
\(E_{21/22}\) route provides an independent computation. &
The recurrence is exact for the declared terminal state and receives neither
\(\alpha\) nor CODATA; the minimal APP deduction\(\to K\) requires construction of its
complete bidirectional reader.\\
\(i\) interno &
\(J^2=-I\) on the Paley--Witt module. &
Construct a finite complex structure and separate conjugate orientations. &
Over the declared base field, \(J^2=-I\) holds exactly.\\
\(-1/12\) &
\(\zeta(-1)=-1/12\) and normalized operator realizations. &
Incidence \(90\to120\), dodecaphase defect and spectral functionals,
each with its domain and normalization. &
Exact equality for each declared map.\\
\(e+\pi\) &
Sum of the presentations of \(e\) and \(\pi\). &
Composite object of the presentation algebra; preserves both genealogies. &
Exact construction; separate arithmetic class.\\
Euler--Mascheroni \(\gamma\) and Stieltjes \(\gamma_n\) &
\(\gamma=\operatorname*{FP}_{s=1}\zeta(s)\); coefficients of Laurent. &
Residual decomposition modulo three and Laurent jets. &
Exact identities of the declared functional.\\
Apéry \(\zeta(3)\) &
Evaluation of the zeta function at \(s=3\). &
Reader \(\mathcal Z_3(3)\) with rational prefix intervals. &
Exact evaluation; classical irrationality.\\
Glaisher--Kinkelin \(A_{\rm GK}\) &
\(A_{\rm GK}=\exp(1/12-\zeta'(-1))\). &
Regularized derivative of the spectral reader. &
Exact identity with explicit normalization.\\
Euler--Kronecker constant \(\gamma_{F_{-3}}\) &
Logarithmic finite part of the Dedekind zeta function of
\(F_{-3}=\mathbb Q(\sqrt{-3})\). &
The character \(\chi_{-3}\) selects the field and organizes
\(L'(1,\chi_{-3})/L(1,\chi_{-3})\). &
Exact identity of the declared character.\\
Meissel--Mertens \(B_1\) &
\(B_1=\lim_{x\to\infty}(\sum_{p\le x}p^{-1}-\log\log x)\). &
Finite part of the trace over primary modes. &
Exact operatorial representation with fixed modes.\\
Catalan \(G\) and Khinchin \(K_0\) &
\(G=L(2,\chi_{-4})\); ergodic mean of the Gauss system. &
Readers of modular character four and continued fraction. &
Classical theorems embedded with their own dynamics.\\
Feigenbaum \((\delta,\alpha_F)\) &
Scales of the renormalization operator in period doubling. &
The dodecaphonic family produces level eight approximants. &
The calculation determines level eight exactly; universality requires
proving convergence to the renormalization fixed point.\\
Rogers--Ramanujan constant \(R(q)\) &
Level-five continued fraction and Klein--Ramanujan identity for
\(j\). &
It organizes level \(5\), \(A_5\), the limit
\(R(q)\to\varphi^{-1}\), and Hecke--Faber replicability. &
Exact modular identities; the comparison of coefficients covers only the
finite depth computed.\\
Prime clock \(p\) &
\(\tau_p=\operatorname{ord}_p(10)\),
\(M_p=(10^{\tau_p}-1)/p\equiv0\pmod9\) for \(p\nmid30\). &
Primes are irreducible periods; composites synchronize clocks
through the least common multiple. &
Exact arithmetic theorem in the declared domain.\\
\bottomrule
\end{longtable}
\endgroup

Finite part, special value, regularized derivative, character, ergodic
average, renormalization scale, and prime period are distinct maps. Their
coexistence within HMT arises from a typed family of operators and
extraction maps.

\subsection{Double projection of a single genealogy}

Let \(\mathfrak X:=X_\infty^{\rm cal}\) be the compact space of enriched
histories with transported initial gauge. Given a declared gauge
\(c\in\mathscr C_{\rm exc}\), two evaluations act on the same
section:
\[
 \mathfrak R:\mathfrak X\longrightarrow\mathbb R,
 \qquad
 \mathfrak E_c:\mathfrak X\longrightarrow\mathbf{Exc},
 \qquad
 \mathfrak E_c(x):=\mathfrak E(x,c),
\]
where \(\mathfrak R\) contrae cylinders and \(\mathfrak E_c\) preserves supports,
incidences, polarity and boundary transport.

\begin{theorem}[Causal unit of the double projection]
\label{p3-83-c44-s2:thm:unidad-causal-doble-proyeccion-numero}
For every history \(\boldsymbol x\in\mathfrak X\), the pair
\[
 \mathfrak D_c(\boldsymbol x)
 :=(\mathfrak R(\boldsymbol x),\mathfrak E_c(\boldsymbol x))
\]
is a joint publication of the same antecedent. The first coordinate
determines the Archimedean value; the second determines the exceptional incidence
class. When the joint readers separate points,
\(\mathfrak D_c\) is injective.
\end{theorem}

\begin{proof}
Both components are defined on the same domain and preserve the
restriction maps. If
\(\mathfrak D_c(\boldsymbol x)=\mathfrak D_c(\boldsymbol y)\), joint separation
of histories implies \(\boldsymbol x=\boldsymbol y\).
\end{proof}

The Archimedean branch produces the cylinders and the values of \(\pi,e,\varphi\). The incidence branch prolongs the same words and the same seal toward Paley, Witt, Golay, the Leech lattice, \(V^\natural\), the Monster group, and Moonshine. The causal relation is the common preimage: both evaluations preserve different invariants of the same history. The real value is the contracted shadow; exceptional symmetry is the incidence memory published by the same state.

\subsection{Genealogical Closure of the Concept of Number}

The concept of number is fixed by the string
\[
 \boxed{
 \text{finite state}
 \longrightarrow\text{compatible extension}
 \longrightarrow\text{genealogical section}
 \longrightarrow
 \begin{cases}
   \text{visible digit},\\
   \text{Archimedean value},\\
   \text{exceptional incidence}.
 \end{cases}}
\]
The digit is a coordinate; the value is a projection; the number is the
complete section; the symmetry is the compatible automorphism that preserves its
memory. Evolutionary conservation transports that identity through the
nonadic return, the helical carry, and the completion
\(729+271=1000\). With these distinctions secured, Dimensional Mechanics
realizes the same genealogy in a new codomain.

\paragraph{Mathematical dependencies.}
The inverse limit supplies the section; the Archimedean reader produces the
value; the binomial completion and the nonadic extension transport unit and
memory; the complementary readers classify the outputs; and the double
projection applies the Archimedean and exceptional maps to the same antecedent.
The Witt--Golay--Leech--Moonshine chain requires, at every step, the
corresponding incidence intertwiner.
